\documentclass[11pt]{article}
\usepackage[a4paper,margin=28mm]{geometry}
\usepackage{amsmath,amssymb,amsthm}
\usepackage[colorlinks=true,linkcolor=blue,urlcolor=blue,citecolor=blue]{hyperref}
\newtheorem*{theorem*}{Theorem}
\newenvironment{statement}[1]{\par\medskip\noindent\phantomsection\textbf{#1}\ }{\par\medskip}
\setlength{\emergencystretch}{3em}
\newcommand{\Q}{\mathbb Q}
\newcommand{\Z}{\mathbb Z}
\newcommand{\F}{\mathbb F}
\newcommand{\tors}{\mathrm{tors}}
\DeclareMathOperator{\End}{End}
\DeclareMathOperator{\Gal}{Gal}
\DeclareMathOperator{\Hom}{Hom}
\DeclareMathOperator{\GL}{GL}
\DeclareMathOperator{\SL}{SL}
\DeclareMathOperator{\tr}{tr}
\DeclareMathOperator{\Frob}{Frob}
\title{Necessary conditions for genus-\(2\) Jacobians split over \(\Q\) with torsion \([25]\), \([32]\), \([2,28]\), \([2,2,20]\) or \([2,2,2,12]\)}
\author{}
\date{15 November 2026}

\begin{document}
\maketitle

\begin{theorem*}
Let \(C/\Q\) be a curve of genus \(2\) whose Jacobian \(J\) is split over \(\Q\). Let \(E\subset J\) be an elliptic
curve over \(\Q\), let \(F=E^\perp\) be its complement for the canonical principal polarization, and let \(n\) be
the gluing degree, so that \(J\cong(E\times F)/\Gamma_\psi\), where \(\Gamma_\psi\) is the graph of a Galois-equivariant
anti-isometry \(\psi\colon E[n]\to F[n]\) and \(n\ge2\). Put \(T_E=E(\Q)_\tors\) and \(T_F=F(\Q)_\tors\), and let \(N_X\)
denote the conductor of \(X\).
\begin{enumerate}
\item[(a)] If \(J(\Q)_\tors\cong\Z/2\times\Z/28\), then \(\{T_E,T_F\}=\{\Z/7,\ \Z/2\times\Z/4\}\), \(n\) is prime to
\(14\), and \(E\) and \(F\) have split multiplicative reduction at \(3\) and at \(5\). In particular
\(v_3(N_J)=v_5(N_J)=2\), so \(225\mid N_J\).
\item[(b)] If \(J(\Q)_\tors\cong(\Z/2)^2\times\Z/20\), then \(\{T_E,T_F\}=\{\Z/10,\ \Z/2\times\Z/4\}\) and \(n\) is
odd and prime to \(5\).
\item[(c)] If \(J(\Q)_\tors\cong(\Z/2)^3\times\Z/12\), then \(\{T_E,T_F\}=\{\Z/2\times\Z/6,\ \Z/2\times\Z/4\}\),
\(n\) is prime to \(6\), and \(E\) and \(F\) have multiplicative reduction at \(2\), \(3\), \(5\), \(7\), \(11\) and
\(13\). In particular \(v_p(N_J)=2\) for these six primes, so \(30030^2=901800900\) divides \(N_J\).
\item[(d)] If \(J(\Q)_\tors\cong\Z/32\), then \(E\) and \(F\) are not isogenous over \(\Q\), and either
\(\{T_E,T_F\}=\{\Z/8,\ \Z/4\}\), \(v_2(n)=2\) and \(j(E)\neq j(F)\), or \(T_E\cong T_F\cong\Z/8\), \(16\mid n\) and
\(E\) and \(F\) are not isogenous over \(\bar\Q\).
\item[(e)] If \(J(\Q)_\tors\cong\Z/25\), then \(T_E\cong T_F\cong\Z/5\), \(5\mid n\), and:
(i) \(E\) and \(F\) are semistable;
(ii) at \(5\), either both have good reduction and \(\{\#\widetilde E(\F_5),\#\widetilde F(\F_5)\}=\{5,10\}\), or both
have split multiplicative reduction, and the second case holds if \(25\mid n\);
(iii) \(66\) divides both \(N_E\) and \(N_F\), or \(210\) divides both;
(iv) relying on the completeness of Cremona's tables below conductor \(500000\), \(N_E\ge8778\), \(N_F\ge8778\) and
\(\max(N_E,N_F)\ge500000\), so that \(N_J=N_EN_F\ge4389000000\).
\end{enumerate}
\end{theorem*}

Here ``split over \(\Q\)'' means that \(J\) is isogenous over \(\Q\) to a product of two elliptic curves over \(\Q\);
then \(J\) contains elliptic curves over \(\Q\), and (I10) below attaches to each of them its complement, the gluing
degree and the anti-isometry. For the Jacobian of a curve of genus \(2\) the gluing degree of every elliptic curve
is at least \(2\). The five groups of the Theorem are the groups of the census that have neither an example nor an
exclusion in the class split over \(\Q\) (15 November 2026):
\href{https://genus-2-torsion.github.io/group.html?g=25}{\([25]\)},
\href{https://genus-2-torsion.github.io/group.html?g=32}{\([32]\)},
\href{https://genus-2-torsion.github.io/group.html?g=2.28}{\([2,28]\)},
\href{https://genus-2-torsion.github.io/group.html?g=2.2.20}{\([2,2,20]\)} and
\href{https://genus-2-torsion.github.io/group.html?g=2.2.2.12}{\([2,2,2,12]\)}. The Theorem excludes none of them:
it lists conditions that an example must satisfy. Part~(d) rests on a statement about all abelian surfaces,
\hyperref[thm-isotypic]{Theorem~5}: an abelian surface over \(\Q\) that is isogenous over \(\Q\) to the square of an
elliptic curve has no rational point of order \(32\); no polarization is needed there.

\emph{Hypotheses.} The condition \(n\ge2\) is needed in (a) and (c): for the product polarization (\(n=1\)) both
statements are false, as 26b1 \(\times\) 15a1 and 30a2 \(\times\) 15a1 show (\hyperref[rem-products]{Remark~5}).
Parts~(b), (d) and~(e) hold for every principally polarized abelian surface over \(\Q\) that is split over \(\Q\),
with the polarization defined over \(\Q\): part~(b) does not use \(n\ge2\), and in (d) and~(e) the torsion forces
\(4\mid n\) and \(5\mid n\). All parts use that the torsion group is exactly the given one; part~(e)(iv) also uses
that \(T_E\) and \(T_F\) are exactly \(\Z/5\) (\hyperref[rem-trace]{Remark~4}).

\emph{Outline.} Section~\ref{gluing} recalls the gluing and determines \(T_E\), \(T_F\) and the gluing degree in
the five cases (\hyperref[prop-parents]{Proposition~3}). In each case a prime \(\ell\mid n\) gives an isomorphism
\(E[\ell]\cong F[\ell]\) of Galois modules. For \([2,28]\) and \([2,2,2,12]\) we compare \(E[\ell]\) and
\(F[\ell]\) at small primes \(p\) (Sections~\ref{sec-22212} and~\ref{sec-228}): the torsion of \(E\) and \(F\)
restricts the reduction types and the traces of Frobenius, and every configuration other than the stated one
gives traces that differ by an integer prime to \(\ell\) or, when \(\ell=p\), different unramified constituents of
\(E[p]\) and \(F[p]\). For \([32]\) (Section~\ref{sec-32}) the curves cannot be isogenous over \(\Q\) by a
lattice argument in the \(2\)-adic Tate module (Section~\ref{sec-isotypic}); two curves with torsion \(\Z/8\)
glued in degree divisible by \(16\) cannot be quadratic twists of each other, by a congruence of traces modulo
\(16\); and for \(\Z/8\) and \(\Z/4\) with the same \(j\)-invariant the intermediate quotient by the \(2\)-torsion
graph is a Weil restriction, whose rational torsion has exponent \(8\). For \([25]\) (Section~\ref{sec-25}) the
point of order \(25\) gives a relation between the Kummer classes of the rational points of order \(5\) of \(E\)
and \(F\) (\hyperref[lemma-kummer]{Lemma~15}), hence a congruence modulo \(25\) between the traces of Frobenius of
\(E\) and \(F\) at all primes, with a scalar that does not depend on the prime. Local arguments at \(5\) and at
small primes then give (i)--(iii); for (iv) we add an obstruction at nonsplit multiplicative primes and a complete
computation in Cremona's tables below conductor \(500000\).

\section{Conventions and inputs}\label{inputs}

All fields have characteristic \(0\), and \(G_k=\Gal(\bar k/k)\). For an abelian variety \(X\) over \(k\) and a
prime \(\ell\), \(T_\ell(X)\) is the Tate module, \(V_\ell(X)=T_\ell(X)\otimes\Q_\ell\) and
\(X[\ell^\infty]=V_\ell(X)/T_\ell(X)\); \(\chi_\ell\colon G_k\to\Z_\ell^\times\) is the \(\ell\)-adic cyclotomic
character and \(\bar\chi_\ell\) its reduction modulo \(\ell\). For an abelian group \(M\), \(M[m]\) is its
\(m\)-torsion and \(M[p^\infty]\) its \(p\)-primary part. We fix embeddings \(\bar\Q\subset\bar\Q_p\), so that
\(G_{\Q_p}\subset G_{\Q}\), and write \(I_p\subset G_{\Q_p}\) for the inertia group; a Frobenius lift is an element
of \(G_{\Q_p}\) inducing \(x\mapsto x^p\) on the residue field. For an elliptic curve \(E\) over \(\Q_p\) with good
reduction, \(\widetilde E\) is its reduction and \(a_p(E)=p+1-\#\widetilde E(\F_p)\); we put \(t_p(E)=a_p(E)\) if
\(E\) has good reduction, and \(t_p(E)=p+1\) or \(t_p(E)=-(p+1)\) if \(E\) has split or nonsplit multiplicative
reduction. \(v_p\) is the \(p\)-adic valuation, \(\Delta\) a discriminant and \(c_4\), \(c_6\) the usual invariants
of a Weierstrass equation.

\begin{itemize}
\item \textbf{(I1) Mazur's theorem} \cite[Theorem~(8)]{Ma}. For an elliptic curve \(E\) over \(\Q\), \(E(\Q)_\tors\)
is \(\Z/m\) with \(1\le m\le10\) or \(m=12\), or \(\Z/2\times\Z/2m\) with \(1\le m\le4\).

\item \textbf{(I2) Olson} \cite{Ol}. An elliptic curve over \(\Q\) with complex multiplication has torsion subgroup
\(0\), \(\Z/2\), \(\Z/3\), \(\Z/4\), \(\Z/6\) or \((\Z/2)^2\); in particular it has no rational point of order \(5\)
or \(8\).

\item \textbf{(I3) Torsion over quadratic fields.} By Kenku--Momose \cite[Theorem~(0.1)]{KM} and Kamienny
\cite{Ka} (see the statement in \cite[Section~1]{Na}), the torsion subgroup of an elliptic curve over a quadratic
field is \(\Z/m\) with \(1\le m\le16\) or \(m=18\), or \(\Z/2\times\Z/2m\) with \(1\le m\le6\), or
\(\Z/3\times\Z/3\), \(\Z/3\times\Z/6\) or \(\Z/4\times\Z/4\). Hence a torsion subgroup containing \(\Z/2\times\Z/8\)
is \(\Z/2\times\Z/8\) (certificate, Part~E).

\item \textbf{(I4) Universal families} \cite[Table~3]{Ku}, also in \cite[Table~3]{HLP}. Put
\(E(b,c)\colon y^2+(1-c)xy-by=x^3-bx^2\). Every pair \((E,P)\) of an elliptic curve over a field \(k\) and a point
\(P\in E(k)\) of order \(5\) is isomorphic to \((E_d,(0,0))\), where \(E_d=E(d,d)\) and \(d\in k^\times\) with
\(d^2-11d-1\neq0\). Every elliptic curve over \(\Q\) whose group of rational points contains a subgroup isomorphic
to \(\Z/2\times\Z/6\) is isomorphic to \(E(c+c^2,c)\) with \(c=(10-2s)/(s^2-9)\) for some \(s\in\Q\).

\item \textbf{(I5) Halving} \cite[Theorem~4.2]{Kn}. If \(E\colon y^2=(x-e_1)(x-e_2)(x-e_3)\) with \(e_i\in k\),
then a point \((x_0,y_0)\in E(k)\) lies in \(2E(k)\) if and only if \(x_0-e_1\), \(x_0-e_2\) and \(x_0-e_3\) are
squares in \(k\).

\item \textbf{(I6) Reduction of elliptic curves} \cite[VII.1--VII.7]{Si1}. Let \(E\) be an elliptic curve over
\(\Q_p\) with a minimal Weierstrass equation (an integral equation with \(v_p(\Delta)<12\) is minimal
\cite[Remark~VII.1.1]{Si1}). The reduction is good, multiplicative or additive according as \(v_p(\Delta)=0\),
\(v_p(\Delta)>0=v_p(c_4)\) or \(v_p(\Delta),v_p(c_4)>0\), and multiplicative reduction is split if the tangent lines
at the node of the reduced curve are defined over \(\F_p\) \cite[Proposition~VII.5.1]{Si1}. There are subgroups
\(E_1(\Q_p)\subseteq E_0(\Q_p)\subseteq E(\Q_p)\) with \(E_1(\Q_p)\) the kernel of reduction, isomorphic to the
group of points of the formal group of \(E\) \cite[Propositions~VII.2.1 and~VII.2.2]{Si1}, and
\(E_0(\Q_p)/E_1(\Q_p)\) is \(\widetilde E(\F_p)\), \(\F_p^\times\), a group of order \(p+1\), or \(\F_p\) in the
good, split, nonsplit and additive cases \cite[Proposition~VII.5.1]{Si1}, \cite[IV.9.4, Step~2]{Si2}. The group
\(E(\Q_p)/E_0(\Q_p)\) has order \(1\) or \(2\) in the nonsplit case \cite[IV.9.4, Step~2]{Si2} and at most \(4\) in
the additive case \cite[Theorem~VII.6.1]{Si1}. The torsion of \(E_1(\Q_p)\) has \(p\)-power order
\cite[Propositions~IV.3.2(b) and~VII.3.1]{Si1}. \(E\) has potentially good reduction if and only if
\(v_p(j(E))\ge0\) \cite[Proposition~VII.5.5]{Si1}; the reduction type does not change in an unramified extension
\cite[Proposition~VII.5.4(a)]{Si1}.

\item \textbf{(I7) The criterion of N\'eron, Ogg and Shafarevich} \cite[Theorem~VII.7.1 and
Corollary~VII.7.2]{Si1}, \cite[\S1, Theorem~1 and Corollary~2]{ST}. Let \(X\) be an abelian variety over a finite
extension \(K\) of \(\Q_p\) and \(\ell\neq p\). Then \(X\) has good reduction if and only if the inertia group of
\(K\) acts trivially on \(T_\ell(X)\), and isogenous abelian varieties over \(K\) have good reduction together.

\item \textbf{(I8) Tate curves} \cite[Theorems~V.3.1 and~V.5.3, Lemma~V.5.2, Corollary~V.5.4 and
Proposition~V.6.1]{Si2}. If
\(E/\Q_p\) has split multiplicative reduction, then \(E\cong E_q\) for a \(q\in p\Z_p\) with
\(v_p(q)=-v_p(j(E))=v_p(\Delta_{\min})\), and \(E(\bar\Q_p)\cong\bar\Q_p^\times/q^{\Z}\) compatibly with
\(G_{\Q_p}\). Hence for \(m\ge1\) there is an exact sequence of \(G_{\Q_p}\)-modules \(0\to\mu_m\to E[m]\to\Z/m\to0\)
in which a point represented by an \(m\)-th root of \(q\) maps to a generator, and similarly
\(0\to\Z_\ell(1)\to T_\ell(E)\to\Z_\ell\to0\). If \(E\) has nonsplit multiplicative reduction, it becomes split
over the unramified quadratic extension of \(\Q_p\), so \(T_\ell(E)\cong\eta\otimes T_\ell(E_q)\) with \(\eta\) the
unramified quadratic character. If \(\ell\ge3\) is a prime not dividing \(v_p(q)\), some element of \(I_p\) acts
on \(E[\ell]\) by a transvection. More generally, if \(v_p(j(E))<0\), then \(E\) is the quadratic twist of the Tate
curve \(E_q\) with \(j(E_q)=j(E)\) by a quadratic character \(\varepsilon\), so \(T_\ell(E)\cong\varepsilon\otimes
T_\ell(E_q)\); \(\varepsilon\) is unramified if and only if \(E\) has multiplicative reduction, by
\cite[Proposition~VII.5.4(a)]{Si1}.

\item \textbf{(I9) Local mod-\(p\) representations} \cite[\S1.11, Propositions~11 and~12, \S1.12,
Proposition~13 and the Corollaries]{Se}. Let \(E\) be an elliptic curve with good ordinary reduction over a finite
extension of \(\Q_p\), of any ramification index. Then the kernel of the reduction map \(E[p]\to
\widetilde E[p](\bar\F_p)\cong\Z/p\) is cyclic of order \(p\), and the map is onto. Now let \(E/\Q_p\) with \(p\ge3\).
If \(E\) has good ordinary reduction, \(I_p\) acts on this kernel by \(\bar\chi_p\) and trivially on the quotient
\(\widetilde E[p](\bar\F_p)\). If \(E\) has good
supersingular reduction, \(I_p\) acts on \(E[p]\) through a cyclic group of order \(p^2-1\), irreducibly. If \(E\)
has multiplicative reduction, \(E[p]\) is as in (I8).

\item \textbf{(I10) Complements and gluing} \cite[Sections~5.3 and~12.1]{BL}, \cite[Section~2]{BD}. Let
\((A,\lambda)\) be a principally polarized abelian surface over \(k\), with \(\lambda\) defined over \(k\), and let
\(E\subset A\) be an elliptic curve over \(k\), with inclusion \(i_E\) and canonical principal polarization
\(\lambda_E\). There is an integer \(n\ge1\), the gluing degree, with \(i_E^\vee\lambda i_E=n\lambda_E\). The
complement \(F=E^\perp\) (the connected component of \(\ker(i_E^\vee\lambda)\)) is an elliptic curve over \(k\) with
\(i_F^\vee\lambda i_F=n\lambda_F\), and \(E\) is the complement of \(F\) \cite[Corollaries~12.1.2
and~12.1.5]{BL}. The homomorphism \(\pi_E=\lambda_E^{-1}i_E^\vee\lambda\colon A\to E\) has kernel \(F\), and
\(\pi_F\colon A\to F\), defined in the same way, has kernel \(E\) \cite[Proposition~12.1.3]{BL}; so \(A/F\cong E\)
and \(A/E\cong F\). Moreover \(\pi_Ei_E=n\), \(\pi_Ei_F=0\), and \(E\cap F=E[n]=F[n]\) inside \(A\)
\cite[Corollary~12.1.4]{BL}. The addition map \(\sigma\colon E\times F\to A\) and \(\alpha=(\pi_E,\pi_F)\colon
A\to E\times F\) satisfy \(\alpha\sigma=[n]\) and \(\sigma\alpha=[n]\) \cite[\S12.1, equation~(3)]{BL}; the kernel
of \(\sigma\) is \(\{(x,-x):x\in E\cap F\}\), the graph \(\Gamma_\psi\) of \(\psi\colon E[n]\to F[n]\),
\(x\mapsto-x\), and \(\sigma^*\lambda=n\lambda_E\times n\lambda_F\) \cite[Lemma~12.1.6]{BL}. Hence \(\Gamma_\psi\)
is isotropic for the product of the Weil pairings, that is, \(\psi\) is an anti-isometry \cite[Lemma~2.4 and the
proof of Lemma~2.6]{BD}. All these maps are defined over \(k\), since the constructions commute with Galois
conjugation; in particular \(\psi\) is \(G_k\)-equivariant. If \(n=1\), then \(\sigma\) is an isomorphism of
principally polarized surfaces \((E\times F,\lambda_E\times\lambda_F)\cong(A,\lambda)\). The Jacobian of a curve of
genus \(2\) is not a polarized product of elliptic curves, since its theta divisor is the irreducible curve
\cite[Corollary~11.8.2]{BL}; so \(n\ge2\) for every elliptic curve in such a Jacobian.

\item \textbf{(I11) Reduction of abelian varieties.} Let \(X\) be an abelian variety with good reduction over the
ring of integers \(O\) of a finite extension \(K\) of \(\Q_p\), with residue field \(\kappa\). The kernel of reduction
\(X(K)\to\widetilde X(\kappa)\) is the group \(\Phi(\mathfrak m)\) of points of a commutative formal group \(\Phi\)
over \(O\) of dimension \(\dim X\), \(\mathfrak m\) the maximal ideal \cite[Theorem~C.2.6]{HS}, and
\(\Phi(\mathfrak m)\) has no torsion of order prime to \(p\) \cite[Proposition~C.2.5]{HS}; so the torsion of
\(X(K)\) of order prime to \(p\) injects into \(\widetilde X(\kappa)\). For an abelian variety over a number field
and a finite place of good reduction, this injectivity is \cite[Theorem~C.1.4]{HS}.

\item \textbf{(I12) Faltings} \cite[\S5, Satz~3, Satz~4 and Korollar~1]{Fal}. For abelian varieties \(X,Y\) over
\(\Q\), \(V_\ell(X)\) is a semisimple \(G_{\Q}\)-module and \(\Hom_{\Q}(X,Y)\otimes\Z_\ell\cong
\Hom_{G_{\Q}}(T_\ell X,T_\ell Y)\). For an elliptic curve \(E_0/\Q\), \(\End_{\Q}(E_0)=\Z\) (complex
multiplications are not defined over \(\Q\)), so \(V_\ell(E_0)\) is irreducible with commutant \(\Q_\ell\), and the
\(\Q_\ell\)-span of the image of \(G_{\Q}\) is \(\End(V_\ell(E_0))\) (density theorem).

\item \textbf{(I13) Conductors} \cite[IV~\S10 and Theorem~IV.10.2]{Si2}, \cite{BK}. For an abelian variety \(X\)
over \(\Q\) and primes \(p\neq\ell\), \(v_p(N_X)=\dim V_\ell(X)-\dim V_\ell(X)^{I_p}+\delta_p(X)\), where the wild
part \(\delta_p(X)\ge0\) vanishes if \(I_p\) acts unipotently on \(V_\ell(X)\). For an elliptic curve,
\(v_p(N_E)\) is \(0\), \(1\) or at least \(2\) according as the reduction is good, multiplicative or additive.

\item \textbf{(I14) Cremona's tables} \cite{Cr}, \cite{LMFDB}. We rely on the completeness of Cremona's tables of
elliptic curves over \(\Q\) of conductor below \(500000\), as recorded by the LMFDB; Magma's
\texttt{CremonaDatabase}, used in Part~H of the certificate, contains these tables (largest conductor \(499999\)).
\end{itemize}

\section{The gluing and the torsion of the two elliptic curves}\label{gluing}

Throughout, \((A,\lambda)\) is a principally polarized abelian surface over \(\Q\), split over \(\Q\), with
\(\lambda\) defined over \(\Q\); \(E\subset A\) is an elliptic curve over \(\Q\), and \(F\), \(n\), \(\psi\),
\(\sigma\), \(\alpha\), \(\pi_E\), \(\pi_F\) are as in (I10). We identify \(T_E\) and \(T_F\) with their images in
\(G=A(\Q)_\tors\).

\begin{statement}{Lemma 1 (the gluing).}\label{lemma-gluing}
\begin{enumerate}
\item[(a)] The sequences \(0\to T_F\to G\to T_E\) and \(0\to T_E\to G\to T_F\), with the maps \(\pi_E\) and
\(\pi_F\), are exact; so \(\#G\) divides \(\#T_E\cdot\#T_F\), and if \(G\) has a point of order \(p\), so does
\(T_E\) or \(T_F\). For a prime \(p\), the \(p\)-rank of \(G\) is at most the sum of those of \(T_E\) and \(T_F\).
\item[(b)] If a prime \(p\) does not divide \(n\), then \(\sigma\) induces an isomorphism of Galois modules
\(E[p^\infty]\times F[p^\infty]\cong A[p^\infty]\), and \(G[p^\infty]\cong T_E[p^\infty]\times T_F[p^\infty]\).
\item[(c)] If \(m\mid n\), then \(\psi\) restricts to an isomorphism of Galois modules \(E[m]\to F[m]\); in
particular \(T_E[m]\cong T_F[m]\), and \(t_p(E)\equiv t_p(F)\pmod m\) at every prime \(p\nmid m\) of good reduction
for \(E\) and \(F\).
\item[(d)] \(V_\ell(A)\cong V_\ell(E)\oplus V_\ell(F)\) as \(G_{\Q}\)-modules for every prime \(\ell\). If \(E\)
and \(F\) have good or multiplicative reduction at \(p\), then \(v_p(N_A)=v_p(N_E)+v_p(N_F)\).
\item[(e)] (Kummer relation.) Let \(R\in A(\Q)\) and write \(R=\sigma(U,V)\) with \(U\in E(\bar\Q)\), \(V\in
F(\bar\Q)\). Then \(\alpha(R)=(nU,nV)\in E(\Q)\times F(\Q)\), and \(gV-V=\psi(gU-U)\) for every \(g\in G_{\Q}\).
\end{enumerate}
\end{statement}

\emph{Proof.} (a) Take rational points in the exact sequences of abelian varieties \(0\to F\to A\to E\to0\) and
\(0\to E\to A\to F\to0\) given by \(\pi_E\) and \(\pi_F\). The \(p\)-rank statement follows from
\(G/T_E\hookrightarrow T_F\). (b) \(\ker\sigma\subseteq(E\times F)[n]\) has order prime to \(p\), so \(\sigma\) is
injective on \((E\times F)[p^\infty]\); it is onto \(A[p^\infty]\), because a preimage of a point of \(A[p^\infty]\)
is a torsion point whose component of order prime to \(p\) maps into \(A[p^\infty]\cap\ker[m]=0\) for some \(m\)
prime to \(p\). Take Galois invariants. (c) \(\psi\) is a \(G_{\Q}\)-equivariant isomorphism \(E[n]\to F[n]\), and
maps \(E[m]=E[n][m]\) onto \(F[m]\); a Frobenius element at \(p\) acts on \(E[m]\) and \(F[m]\) with traces
\(a_p(E)\) and \(a_p(F)\) modulo \(m\). (d) \(\sigma\) is an isogeny over \(\Q\). By (I13), \(\delta_p=0\) for
\(E\), \(F\) and \(A\) when \(I_p\) acts unipotently on \(V_\ell(E)\) and \(V_\ell(F)\) (\(\ell\neq p\)), which
holds for good and multiplicative reduction by (I7) and (I8), and the term \(\dim V_\ell-\dim V_\ell^{I_p}\) is
additive. (e) \(\alpha(R)=\alpha\sigma(U,V)=(nU,nV)\), which is rational since \(\alpha\) is defined over \(\Q\). For
\(g\in G_{\Q}\), \(gR=R\) gives \(\sigma(gU-U,gV-V)=0\), so \((gU-U,gV-V)\in\ker\sigma=\Gamma_\psi\). \qed

\begin{statement}{Proposition 3 (the torsion of the two curves).}\label{prop-parents}
Let \(G=A(\Q)_\tors\). After exchanging \(E\) and \(F\) if necessary:
\begin{enumerate}
\item[(a)] if \(G\cong\Z/25\), then \(T_E\cong T_F\cong\Z/5\) and \(5\mid n\);
\item[(b)] if \(G\cong\Z/32\), then either \(T_E\cong\Z/8\), \(T_F\cong\Z/4\) and \(v_2(n)=2\), or
\(T_E\cong T_F\cong\Z/8\) and \(16\mid n\);
\item[(c)] if \(G\cong\Z/2\times\Z/28\), then \(T_E\cong\Z/7\), \(T_F\cong\Z/2\times\Z/4\) and \(n\) is prime to
\(14\);
\item[(d)] if \(G\cong(\Z/2)^2\times\Z/20\), then \(T_E\cong\Z/10\), \(T_F\cong\Z/2\times\Z/4\) and \(n\) is odd
and prime to \(5\);
\item[(e)] if \(G\cong(\Z/2)^3\times\Z/12\), then \(T_E\cong\Z/2\times\Z/6\), \(T_F\cong\Z/2\times\Z/4\) and \(n\)
is prime to \(6\).
\end{enumerate}
If \(A\) is the Jacobian of a curve of genus \(2\), then \(n\ge3\) in (c) and (d), and \(n\ge5\) in (e).
\end{statement}

\emph{Proof.} By \hyperref[lemma-gluing]{Lemma~1}(a), \(T_E\) and \(T_F\) are subgroups of \(G\) from the list
(I1), \(\#G\mid\#T_E\#T_F\), and \(G/T_E\) and \(G/T_F\) embed in \(T_F\) and \(T_E\). Part~A of the certificate
lists, for each of the five groups \(G\), all pairs of groups from (I1) with these properties; they are the pairs
in (a)--(e). By hand:
\begin{itemize}
\item \(G=\Z/25\): \(T_E,T_F\subseteq\Z/25\) are \(0\) or \(\Z/5\) by (I1), and \(25\mid\#T_E\#T_F\).
\item \(G=\Z/32\): \(T_E\), \(T_F\) are cyclic \(2\)-groups, so of order at most \(8\), with product of orders at
least \(32\).
\item \(G=\Z/2\times\Z/28\): one of the curves, say \(E\), has a point of order \(7\), so \(T_E=\Z/7\), of
\(2\)-rank \(0\); then \(T_F\) has \(2\)-rank \(2\) and embeds in \(G\), so \(T_F\in\{(\Z/2)^2,\Z/2\times\Z/4\}\),
and \(56\mid7\#T_F\) gives \(T_F=\Z/2\times\Z/4\).
\item \(G=(\Z/2)^2\times\Z/20\): one curve, say \(E\), has a point of order \(5\), so \(T_E\in\{\Z/5,\Z/10\}\) has
\(2\)-rank at most \(1\); as \(G\) has \(2\)-rank \(3\), \(T_E=\Z/10\) and \(T_F\) has \(2\)-rank \(2\); it embeds in
\(G\), whose exponent is \(20\), and \(80\mid10\#T_F\), so \(T_F=\Z/2\times\Z/4\).
\item \(G=(\Z/2)^3\times\Z/12\): both have \(2\)-rank \(2\), and one, say \(E\), has a point of order \(3\), so
\(T_E=\Z/2\times\Z/6\). Then \(T_F=\Z/2\times\Z/2m\) embeds in \(G\), so \(m\le3\), and \(\#T_F\ge96/12=8\). If
\(T_F=\Z/2\times\Z/6\), then \(G/T_E\), of order \(8\), embeds in \(T_F\), whose \(2\)-part has order \(4\), which
is absurd. So \(T_F=\Z/2\times\Z/4\).
\end{itemize}
Gluing degrees. If a prime \(p\) divides \(n\), then \(T_E[p]\cong T_F[p]\) by
\hyperref[lemma-gluing]{Lemma~1}(c). In (c) this excludes \(p=2\) and \(p=7\); in (d) \(p=2\) and \(p=5\); in (e)
\(p=3\).

(a) Let \(R\) generate \(G\). Since \(T_E\) and \(T_F\) are killed by \(5\), \(5R\) lies in \(\ker\pi_E=F\) and in
\(\ker\pi_F=E\), so \(5R\in E\cap F=E[n]\) is a point of order \(5\), and \(5\mid n\).

(b) Both \(T_E\) and \(T_F\) have order at least \(4\), so both contain the unique subgroup of order \(4\) of the
cyclic group \(G\), which therefore lies in \(E\cap F=E[n]\); so \(4\mid n\). If \(8\mid n\), then
\(T_E[8]\cong T_F[8]\), so \(T_E\cong T_F\cong\Z/8\). If \(T_E\cong T_F\cong\Z/8\), both equal the subgroup \(H\) of
order \(8\) of \(G\), and \(H\subseteq E\cap F=E[n]\) gives \(8\mid n\). Suppose \(v_2(n)=3\) and let \(R\) generate
\(G\). The kernel of \(\pi_E\) on \(G\) is \(T_F=H\), so \(\pi_E(R)\in T_E=H\) has order \(4\) and equals \(8uR\)
with \(u\) odd; likewise \(\pi_F(R)=8vR\) with \(v\) odd. Then \(nR=\sigma\alpha(R)=8(u+v)R\) has order at most
\(2\), whereas \(n=8m\) with \(m\) odd and \(8mR\) has order \(4\). So \(16\mid n\).

(e) It remains to show that \(n\) is odd. Suppose not, let \(q=2^{v_2(n)}\ge2\), and let \(\psi_q\colon E[q]\to
F[q]\) be the restriction of \(\psi\). Choose \(\Z_2\)-bases \((e_1,e_2)\) of \(T_2(E)\) and \((f_1,f_2)\) of
\(T_2(F)\) with \(\psi_q(e_i/q)=f_i/q\) (lift \(\psi_q(e_i/q)\) and use Nakayama's lemma), where \(v/q\) denotes the
image of \(v\in T_2\) in \(T_2/qT_2=X[q]\). Let \(M_E(g),M_F(g)\in\GL_2(\Z_2)\) be the matrices of \(g\in G_{\Q}\).
Equivariance of \(\psi_q\) gives \(M_E(g)\equiv M_F(g)\pmod q\); put \(D(g)=(M_F(g)-M_E(g))/q\). The \(2\)-primary part
of \(\ker\sigma\) is the graph of \(\psi_q\), and
\[\theta\colon(\Q_2/\Z_2)^2\times(\Q_2/\Z_2)^2\to A[2^\infty],\qquad
\theta(x,z)=\sigma\bigl(\tilde x/q,\ \tilde x/q+\tilde z\bigr),\]
where \(\tilde x,\tilde z\in\Q_2^2\) are lifts and a vector \(w\in\Q_2^2\) stands for the point \(w_1e_1+w_2e_2\) of
\(E[2^\infty]\), resp.~\(w_1f_1+w_2f_2\) of \(F[2^\infty]\), is a well-defined isomorphism: changing \(\tilde x\) by
\(u\in\Z_2^2\) changes the pair by \((u/q,u/q)\in\Gamma_{\psi_q}\); \(\theta(x,z)=0\) forces \(\tilde x\in\Z_2^2\)
and then \(\tilde z\in\Z_2^2\); and \(\sigma(s,t)=\theta(qs,t-s)\). It satisfies
\(g\,\theta(x,z)=\theta\bigl(M_E(g)x,\ D(g)x+M_F(g)z\bigr)\); this is the model of \cite[Section~2]{P256}. Since \(G\) has \(2\)-rank \(4\), \(A[2]\subseteq
A(\Q)\), so \(G_{\Q}\) acts trivially on \(\theta\bigl((\tfrac12\Z_2^2/\Z_2^2)^2\bigr)=A[2]\): \(M_E(g)\equiv
M_F(g)\equiv I\) and \(D(g)\equiv0\pmod 2\). Hence \(M_F(g)-M_E(g)=qD(g)\equiv0\pmod{2q}\), and the identity matrix
gives an isomorphism of Galois modules \(E[2q]\to F[2q]\). As \(2q\ge4\), \(T_E[4]\cong T_F[4]\); but these groups
are \((\Z/2)^2\) and \(\Z/2\times\Z/4\).

Finally, if \(A\) is a Jacobian then \(n\ge2\) by (I10); in (c) and (d) \(n\) is odd, and in (e) prime to \(6\).
\qed

\medskip
For \([2,2,20]\) we prove nothing beyond \hyperref[prop-parents]{Proposition~3}(d). The condition \(n\ge2\) enters
\hyperref[prop-parents]{Proposition~3} only in its last sentence.

\section{Surfaces isogenous to a square}\label{sec-isotypic}

\begin{statement}{Lemma 4 (stable lattices).}\label{lemma-lattices}
Let \(E_0/\Q\) be an elliptic curve, \(p\) a prime, \(V=V_p(E_0)\), \(c\in G_{\Q}\) a complex conjugation and
\(s_{+},s_{-}\in V\) eigenvectors of \(c\) with eigenvalues \(1\) and \(-1\). Let \(W\) be a finite-dimensional
\(\Q_p\)-vector space with trivial \(G_{\Q}\)-action, and \(T\subset V\otimes W\) a \(G_{\Q}\)-stable
\(\Z_p\)-lattice of the form \(T=(s_{+}\otimes\Lambda_{+})\oplus(s_{-}\otimes\Lambda_{-})\) for lattices
\(\Lambda_\pm\subset W\). Then \(T\cong\bigoplus_jT_p(E_j)\) as \(\Z_p[G_{\Q}]\)-modules, for elliptic curves
\(E_1,\dots,E_r\) over \(\Q\) that are isogenous over \(\Q\) to \(E_0\), \(r=\dim W\).
\end{statement}

\emph{Proof.} By elementary divisors there are a basis \(w_1,\dots,w_r\) of \(W\) and integers \(d_j\) with
\(\Lambda_{+}=\bigoplus\Z_pw_j\) and \(\Lambda_{-}=\bigoplus p^{d_j}\Z_pw_j\). For \(g\in G_{\Q}\) write
\(gs_{+}=as_{+}+c's_{-}\) and \(gs_{-}=bs_{+}+ds_{-}\) with \(a,b,c',d\in\Q_p\). Since \(g(s_{+}\otimes w_j)\) and
\(g(p^{d_j}s_{-}\otimes w_j)\) lie in \(T\), we get \(a,d\in\Z_p\), \(c'\in p^{d_j}\Z_p\) and \(p^{d_j}b\in\Z_p\). So
\(S_j=\Z_p\,s_{+}\otimes w_j\oplus p^{d_j}\Z_p\,s_{-}\otimes w_j\) is \(G_{\Q}\)-stable and \(T=\bigoplus S_j\). Each
\(S_j\) is a \(G_{\Q}\)-stable lattice in \(V\otimes w_j\cong V\); after scaling it contains \(T_p(E_0)\), and then
\(S_j/T_p(E_0)\) is a finite Galois-stable subgroup \(C_j\) of \(E_0[p^\infty]\), and \(S_j\cong T_p(E_0/C_j)\).
\qed

\begin{statement}{Theorem 5 (surfaces isogenous to a square).}\label{thm-isotypic}
Let \(E_0/\Q\) be an elliptic curve and \(B/\Q\) an abelian surface that is isogenous over \(\Q\) to
\(E_0\times E_0\).
\begin{enumerate}
\item[(a)] For every odd prime \(p\) there are elliptic curves \(E_1,E_2\) over \(\Q\), isogenous over \(\Q\) to
\(E_0\), with \(B(\Q)[p^\infty]\cong E_1(\Q)[p^\infty]\times E_2(\Q)[p^\infty]\). In particular \(B(\Q)\) has no
point of order \(25\).
\item[(b)] The exponent of \(B(\Q)[2^\infty]\) divides \(16\). In particular \(B(\Q)\) has no point of order
\(32\).
\end{enumerate}
\end{statement}

\emph{Proof.} An isogeny \(E_0^2\to B\) over \(\Q\) identifies \(V_p(B)\) with \(V\otimes W\), \(V=V_p(E_0)\),
\(W=\Q_p^2\) with trivial action; let \(T\) be the image of \(T_p(B)\). Complex conjugation \(c\) has
\(\det c=\chi_p(c)=-1\) and \(c^2=1\), so it has eigenvalues \(1\) and \(-1\) on \(V\), and \(e_\pm=(1\pm c)/2\)
act on \(V\otimes W\) through the first factor.

(a) For odd \(p\), \(e_\pm\in\Z_p[G_{\Q}]\) preserve \(T\), so \(T=e_{+}T\oplus e_{-}T\) has the form of
\hyperref[lemma-lattices]{Lemma~4}, and \(B[p^\infty]=V_p(B)/T\cong E_1[p^\infty]\times E_2[p^\infty]\) as Galois
modules. By (I1), \(E_i(\Q)[5^\infty]\) has exponent at most \(5\).

(b) If no elliptic curve isogenous over \(\Q\) to \(E_0\) has a rational point of order \(4\), let \(E'\subset B\)
be the image of \(E_0\times0\); then \(E'\) and \(B/E'\) are elliptic curves over \(\Q\) isogenous to \(E_0\) (by
(I12), since \(V_2(B/E')\cong V_2(B)/V_2(E')\cong V_2(E_0)\)), their rational \(2\)-primary torsion has exponent at
most \(2\), and the exact sequence \(0\to E'(\Q)\to B(\Q)\to(B/E')(\Q)\) shows that \(B(\Q)[2^\infty]\) has
exponent at most \(4\). Otherwise we may replace \(E_0\) by an isogenous curve with a rational point of order \(4\).

\emph{Step 1: a stable lattice on which \(c\) is diagonal.} Let \(e_1\) be a primitive vector of \(T_2(E_0)\) with
\(ce_1=e_1\), and complete it to a basis \((e_1,e_2)\); then \(c=\left(\begin{smallmatrix}1&b\\0&-1
\end{smallmatrix}\right)\) (columns are images of basis vectors), and replacing \(e_2\) by \(e_2+ke_1\) changes \(b\)
by \(2k\), so we may take \(b\in\{0,1\}\). If \(b=0\), put \(S=T_2(E_0)\). If \(b=1\), a rational point of order
\(4\) is \((v_1e_1+v_2e_2)/4\) with \((v_1,v_2)\) primitive modulo \(4\) and fixed by \(c\); the primitive vectors
modulo \(4\) fixed by \(c\) are the odd multiples of \(e_1\) (certificate, Part~E), so \(e_1/4\) is rational, and
every \(g\in G_{\Q}\) has a matrix \(\left(\begin{smallmatrix}a&b_0\\c_0&d\end{smallmatrix}\right)\) with
\(a\equiv1\), \(c_0\equiv0\pmod4\) and \(d\) odd. In the basis \(s_{+}=e_1\), \(s_{-}=e_2-e_1/2\) of \(V\), \(c\) is
diagonal and \(g\) has the matrix
\[\begin{pmatrix}a+c_0/2&b_0+(d-a)/2-c_0/4\\ c_0&d-c_0/2\end{pmatrix},\]
which is integral (certificate, Part~E). So \(S=\Z_2s_{+}\oplus\Z_2s_{-}\) is \(G_{\Q}\)-stable, with \(cs_\pm=\pm
s_\pm\).

\emph{Step 2: an overorder.} Let \(R\) be the \(\Z_2\)-span of the image of \(G_{\Q}\) in \(\End(S)\cong M_2(\Z_2)\)
(matrices in the basis \(s_\pm\)). By (I12), \(R\otimes\Q=M_2(\Q_2)\), so the sets \(I_{12}\), \(I_{21}\) of
upper-right and lower-left entries of elements of \(R\) are nonzero ideals of \(\Z_2\), and
\(R_0=\left(\begin{smallmatrix}\Z_2&I_{12}\\ I_{21}&\Z_2\end{smallmatrix}\right)\) is an order containing \(R\)
and \(e_\pm\). Moreover \(2R_0\subseteq R\): \(2e_{+}=1+c\) and \(2e_{-}=1-c\) lie in \(R\), and for \(r\in R\) with
diagonal entries \(r_{11}\), \(r_{22}\),
\[(1+c)r-r_{11}(1+c)=2e_{+}re_{-},\qquad (1-c)r-r_{22}(1-c)=2e_{-}re_{+}\]
(certificate, Part~E), and these are \(2r_{12}\) and \(2r_{21}\) placed in the corners.

\emph{Step 3.} \(T'=R_0T\) is \(G_{\Q}\)-stable and \(e_\pm\)-stable, and \(T\subseteq T'\subseteq\tfrac12T\) since
\(2R_0\subseteq R\) and \(RT=T\). By \hyperref[lemma-lattices]{Lemma~4} (with \(V=\Q_2S\)), \(T'\cong
T_2(E_1)\oplus T_2(E_2)\), so \((V_2(B)/T')^{G_{\Q}}\cong E_1(\Q)[2^\infty]\times E_2(\Q)[2^\infty]\) has exponent
dividing \(8\) by (I1). For \(x\in B(\Q)[2^\infty]=(V_2(B)/T)^{G_{\Q}}\), the image of \(x\) in \(V_2(B)/T'\) is
rational, so \(8\tilde x\in T'\) for a lift \(\tilde x\), and \(16\tilde x\in2T'\subseteq T\). \qed

\medskip
No polarization is used, and (b) also holds when \(E_0\) has complex multiplication.

\section{Local tools}\label{local}

\begin{statement}{Lemma 6 (formal groups).}\label{lemma-formal}
Let \(O\) be a complete discrete valuation ring of characteristic \(0\) with residue characteristic \(p\),
normalized valuation \(v\) and \(e=v(p)\), and let \(\Phi\) be a commutative formal group law of dimension \(g\) over
\(O\). For \(x\in\mathfrak m^g\) put \(v(x)=\min_iv(x_i)\).
\begin{enumerate}
\item[(a)] Every coefficient of \([p](X)\) is in \(pO\), except those of the monomials \(X^\gamma\) with
\(\gamma\in p\mathbb N^g\smallsetminus\{0\}\). Hence \(v([p](x))\ge\min(e+v(x),\ pv(x))\).
\item[(b)] If \(x\neq0\) and \([p](x)=0\), then \(v(x)\le e/(p-1)\).
\item[(c)] If \(e<p-1\), then \(\Phi(\mathfrak m)\) has no point of order \(p\); if \(e<p(p-1)\), it has no point of
order \(p^2\).
\item[(d)] For an elliptic curve \(E/\Q_p\), \(E_1(\Q_p)\) is torsion-free if \(p\ge3\), and its torsion has order
at most \(2\) if \(p=2\).
\end{enumerate}
\end{statement}

\emph{Proof.} (a) Let \(\Phi_p(X_1,\dots,X_p)=\sum_\beta a_\beta X_1^{\beta_1}\cdots X_p^{\beta_p}\) be the
\(p\)-fold sum; it has coefficients in \(O\) and is symmetric, by commutativity and associativity. The coefficient
of \(X^\gamma\) in \([p](X)=\Phi_p(X,\dots,X)\) is \(\sum a_\beta\) over the \(\beta\) with
\(\beta_1+\dots+\beta_p=\gamma\). The cyclic permutations of the \(p\) factors preserve \(a_\beta\), and their orbits
have \(p\) elements except for the fixed points \(\beta_1=\dots=\beta_p\), which occur only if \(\gamma\in
p\mathbb N^g\). So the coefficient lies in \(pO\) unless \(\gamma\in p\mathbb N^g\). A monomial with coefficient in
\(pO\) and total degree at least \(1\) has valuation at least \(e+v(x)\) at \(x\), and one with \(\gamma\in
p\mathbb N^g\smallsetminus\{0\}\) has valuation at least \(pv(x)\).

(b) Let \(s=v(x)\) and suppose \(s>e/(p-1)\); choose \(i\) with \(v(x_i)=s\). The \(i\)-th coordinate of \([p](x)\)
is \(px_i\), of valuation \(e+s\), plus monomials of valuation at least \(e+2s\) (coefficient in \(pO\), degree at
least \(2\)) or \(ps>e+s\). So it is nonzero.

(c) The first claim follows from (b), since \(v(x)\ge1\). If \(x\) has order \(p^2\), then \(y=[p](x)\ne0\) is
killed by \(p\), and by (a), (b): \(e/(p-1)\ge v(y)\ge\min(e+1,p)\), so \(e\ge p(p-1)\).

(d) \(E_1(\Q_p)\) is the group of points of the formal group of \(E\) over \(\Z_p\) (I6), and \(e=1\). For \(p\ge3\)
use (c) and the fact that its torsion has \(p\)-power order (I6). For \(p=2\), the points with \(v\ge2\) form a
torsion-free group, isomorphic to \(4\Z_2\) by the formal logarithm \cite[Theorem~IV.6.4(b)]{Si1}, of index \(2\)
\cite[Proposition~IV.3.2(a)]{Si1}. \qed

\begin{statement}{Lemma 7 (local torsion).}\label{lemma-localtorsion}
Let \(E/\Q_p\) be an elliptic curve.
\begin{enumerate}
\item[(a)] If \(E\) has additive reduction, \(E(\Q_p)_\tors\) has order \(ap^k\) with \(a\le4\).
\item[(b)] If \(E\) has nonsplit multiplicative reduction, the part of \(E(\Q_p)_\tors\) of order prime to \(p\) has
order dividing \(2(p+1)\); if \(p\ge3\), so does all of \(E(\Q_p)_\tors\); if \(p=2\), \(E(\Q_2)[2^\infty]\) has
order at most \(4\).
\item[(c)] If \(E\) has good reduction, \(E(\Q_p)_\tors\) injects into \(\widetilde E(\F_p)\) if \(p\ge3\), and the
kernel of \(E(\Q_2)_\tors\to\widetilde E(\F_2)\) has order at most \(2\).
\end{enumerate}
\end{statement}

\emph{Proof.} Use the filtration of (I6) and \hyperref[lemma-formal]{Lemma~6}(d). (a) \(E_1(\Q_p)\) and
\(E_0/E_1\cong\F_p\) contribute a power of \(p\) and \(E/E_0\) at most \(4\) (\cite[Remark~IV.9.2.2]{Si2}). (b)
\(E_0/E_1\) has order \(p+1\) and \(E/E_0\) order at most \(2\); for \(p=2\), \(E_1(\Q_2)\) contributes at most
\(2\), \(E_0/E_1\) (of order \(3\)) nothing and \(E/E_0\) at most \(2\). (c) The kernel of reduction is \(E_1\). \qed

\begin{statement}{Lemma 8 (comparison of local representations).}\label{lemma-compare}
Let \(p\) be a prime, \(\ell\ge3\) a prime, and \(E\), \(F\) elliptic curves over \(\Q_p\) with good or
multiplicative reduction such that \(E[\ell]\cong F[\ell]\) as \(G_{\Q_p}\)-modules.
\begin{enumerate}
\item[(a)] If \(\ell\neq p\), then \(t_p(E)\equiv t_p(F)\pmod\ell\).
\item[(b)] If \(\ell=p\), then \(E\) has good supersingular reduction if and only if \(F\) has. Otherwise
\(u(E)=u(F)\) in \(\F_p\), where \(u(X)=a_p(X)\bmod p\) if \(X\) has good ordinary reduction and \(u(X)=\pm1\) if
\(X\) has split, resp.~nonsplit, multiplicative reduction.
\end{enumerate}
\end{statement}

\emph{Proof.} (a) Let \(\varphi\) be a Frobenius lift. If \(X\in\{E,F\}\) has good reduction, \(X[\ell]\) is
unramified and \(\varphi\) has trace \(a_p(X)\) on it; if \(X\) has multiplicative reduction, the semisimplification
of \(X[\ell]\) is \(\varepsilon\oplus\varepsilon\bar\chi_\ell\) with \(\varepsilon\) trivial or \(\eta\) (I8), and
\(\varphi\) has trace \(\varepsilon(\varphi)(1+p)=t_p(X)\). Isomorphic modules give equal traces. (b) By (I9), if
\(X\) is supersingular then \(X[p]\) is irreducible, and otherwise the semisimplification of \(X[p]\) is
\(\psi_X\oplus\bar\chi_p\psi_X^{-1}\) with \(\psi_X\) unramified, \(\bar\chi_p\psi_X^{-1}\) ramified (as
\(\bar\chi_p\) has order \(p-1\ge2\) on \(I_p\)), and \(\psi_X(\varphi)=u(X)\). In the multiplicative case this is
(I8). In the ordinary case, \(\psi_X\) is the action on \(\widetilde X[p](\bar\F_p)\), on which \(\varphi\) acts as
the Frobenius endomorphism \(\pi\); since \(\pi\) is bijective on points and \(\hat\pi\pi=p\), the dual \(\hat\pi\)
kills \(\widetilde X[p](\bar\F_p)=\pi(\widetilde X[p](\bar\F_p))\), so \(\pi=a_p-\hat\pi\) acts as \(a_p\).
Isomorphic modules have the same semisimplification, hence the same unramified constituent. \qed

\section{The group \([2,2,2,12]\)}\label{sec-22212}

Let \(A_t\), \(B_t\) and the curves \(E_t\), \(F_t\) over \(\Q(t)\) be
\[A_t=(t+3)(t-1)^3,\qquad B_t=(t-3)(t+1)^3,\]
\[E_t\colon y^2=x(x-A_t)(x-B_t),\qquad F_t\colon y^2=x(x+1)(x+t^2).\]
The curve \(E_t\) is the model \(F_{2,6}^t\) of Howe, Lepr\'evost and Poonen \cite[Tables~6 and~8]{HLP}.

\begin{statement}{Lemma 9 (universal families and valuations).}\label{lemma-hlp}
\begin{enumerate}
\item[(a)] \(E_t\) has full rational \(2\)-torsion and the point \((t^4-2t^2+1,4(t^2-1)^2)\) of order \(3\) over
\(\Q(t)\), and
\[j(E_t)=\frac{256(A_t^2-A_tB_t+B_t^2)^3}{A_t^2B_t^2(A_t-B_t)^2},\qquad A_t-B_t=16t,\]
of degree \(24\) in \(t\); \(E_t\) is nonsingular exactly for \(t\notin\{0,\pm1,\pm3\}\), and the poles of
\(j(E_t)\) are \(t\in\{0,\pm1,\pm3,\infty\}\). Every elliptic curve
\(E\) over \(\Q\) with \(\Z/2\times\Z/6\subseteq E(\Q)\) satisfies \(j(E)=j(E_t)\) for some \(t\in\Q\smallsetminus
\{0,\pm1,\pm3\}\).
\item[(b)] \(F_t\) has full rational \(2\)-torsion and the point \((t,t(t+1))\) of order \(4\), and
\(j(F_t)=256(t^4-t^2+1)^3/(t^4(t^2-1)^2)\). Every elliptic curve \(F\) over \(\Q\) with \(\Z/2\times\Z/4\subseteq
F(\Q)\) is isomorphic over \(\Q\) to \(F_t\) for some \(t\in\Q\smallsetminus\{0,\pm1\}\).
\item[(c)] For every \(t\in\Q\smallsetminus\{0,\pm1,\pm3\}\), \(v_2(j(E_t))<0\) and \(v_3(j(E_t))<0\); for every
\(t\in\Q\smallsetminus\{0,\pm1\}\), \(v_3(j(F_t))<0\).
\end{enumerate}
\end{statement}

\emph{Proof.} The identities and degrees in (a) and (b) are checked in Part~B of the certificate. (a) By (I4),
\(E\cong E(c+c^2,c)\) with \(c=(10-2s)/(s^2-9)\), and the certificate checks that the \(j\)-invariant of this family is
\(j(E_{(s-3)/2})\); as \(j(E)\) is finite, \(t=(s-3)/2\) is not a pole of \(j(E_t)\), so
\(t\notin\{0,\pm1,\pm3\}\).
Alternatively, \(E_t\) carries a level structure \(\Z/2\times\Z/6\) over \(\Q(t)\), so it defines
a map from the \(t\)-line to the modular curve \(X_1(2,6)\), the coarse moduli space of curves with such a level
structure \cite[p.~125]{KM}; the map \(X_1(2,6)\to X(1)\) has degree \([\SL_2(\Z):\{\pm1\}\Gamma]=24\)
\cite[Exercise~3.1.1]{DS}, where \(\Gamma\) is the group of matrices fixing \(e_1\) modulo \(6\) and \(e_2\) modulo
\(2\) (certificate, Part~B), so the map from the \(t\)-line has degree \(24/24=1\) and is an isomorphism. (b) Let
\(P\in F(\Q)\) have order \(4\) and move \(T=2P\) to \((0,0)\): \(F\colon y^2=x(x-e_2)(x-e_3)\) with \(e_2,e_3\in\Q\).
By (I5), \(-e_2\) and \(-e_3\) are squares, \(-e_2=r^2\), \(-e_3=r^2t^2\), and \((x,y)\mapsto(x/r^2,y/r^3)\) gives
\(F_t\); \(t\neq0,\pm1\) as \(F\) is nonsingular. (c) For \(F_t\) write \(t=a/b\) in lowest terms:
\(j(F_t)=256(a^4-a^2b^2+b^4)^3/(a^4b^4(a^2-b^2)^2)\). If \(3\nmid ab\), the numerator is \(\equiv256\not\equiv0\)
and \(3\mid a^2-b^2\); if \(3\mid ab\), the numerator is prime to \(3\) and the denominator is not. For \(E_t\), with
\(C=a^4-6a^2b^2-3b^4\) and \(D=8ab^3\), one has \(A=C+D\), \(B=C-D\) and \(A^2-AB+B^2=C^2+3D^2\) for the homogenized
\(A_t\), \(B_t\) (certificate). At \(2\): if \(a\), \(b\) have opposite parity, \(A\), \(B\) and \(C^2+3D^2\) are odd
and \(v_2(j)=8-2v_2(16ab^3)<0\); if both are odd, one of \(A\), \(B\) has valuation \(4\) and the other valuation
\(r\ge9\), and \(v_2(j)=8+24-8-2r-8<0\). At \(3\): if \(3\nmid a\), the numerator is a unit and \(3\) divides
\((a-b)(a+b)b\), hence the denominator; if \(k=v_3(a)\ge1\), then \(v_3(C)=1\), \(v_3(C^2+3D^2)=2\), and
\(v_3(j)=6-(4+2k)<0\) for \(k\ge2\), while for \(k=1\) one of \(A\), \(B\) has valuation \(1\), the other at least
\(2\), and \(v_3(j)\le6-8<0\). Part~B of the certificate also proves (c) mechanically on all of \(\mathbb
P^1(\Q_2)\) and \(\mathbb P^1(\Q_3)\): it covers them by finitely many discs on each of which \(v_p\) of the
numerator is constant and smaller than \(v_p\) of the denominator. \qed

\begin{statement}{Theorem 10 (the group [2,2,2,12]).}\label{thm-22212}
Let \(A\), \(E\), \(F\), \(n\) be as in Section~\ref{gluing} with \(A(\Q)_\tors\cong(\Z/2)^3\times\Z/12\) and
\(n\ge2\). Then \(E\) and \(F\) have multiplicative reduction at \(2\), \(3\), \(5\), \(7\), \(11\) and \(13\); they
are split at \(2\), \(3\), \(5\), \(7\), \(13\), and have the same splitting type at \(11\). Hence \(v_p(N_A)=2\) for
these primes and \(901800900=30030^2\) divides \(N_A\).
\end{statement}

\emph{Proof.} By \hyperref[prop-parents]{Proposition~3}(e), \(T_E=\Z/2\times\Z/6\), \(T_F=\Z/2\times\Z/4\), and \(n\)
is prime to \(6\); as \(n\ge2\), it has a prime factor \(\ell\ge5\), and \(E[\ell]\cong F[\ell]\) by
\hyperref[lemma-gluing]{Lemma~1}(c). We apply \hyperref[lemma-compare]{Lemma~8} to this \(\ell\) at each \(p\), once
we know that both curves are semistable at \(p\).

\emph{The primes \(2\) and \(3\).} By \hyperref[lemma-hlp]{Lemma~9}(a) and~(c), \(v_2(j(E))<0\) and
\(v_3(j(E))<0\), and by Lemma~9(b) and~(c), \(v_3(j(F))<0\); so these reductions are multiplicative or additive (I6).
If \(F\) were additive at \(3\), its \(8\) rational \(2\)-primary points would contradict
\hyperref[lemma-localtorsion]{Lemma~7}(a); so \(F\) is multiplicative at \(3\). If \(E\) were additive at \(2\),
then, as \(v_2(j(E))<0\), \(E\) would be the twist of a Tate curve by a ramified quadratic character \(\varepsilon\)
(I8); an element \(\sigma\in I_2\) with \(\varepsilon(\sigma)=-1\) acts on \(E[3]\) as \(-1\) times a unipotent
matrix, since \(\chi_3\) is trivial on \(I_2\), so it fixes no nonzero point of \(E[3]\), and \(E(\Q_2)\) would have
no point of order \(3\). So \(E\) is multiplicative at \(2\). If \(E\) were additive at \(3\),
then \(v_3(j(E))<0\) makes \(E\) a ramified quadratic twist of a Tate curve (I8), so some element of \(I_3\) acts on
\(E[\ell]\) with both eigenvalues \(-1\), while \(I_3\) acts unipotently on \(F[\ell]\); so \(E\) is multiplicative
at \(3\). At \(2\), \(F\) is not additive: in the potentially multiplicative case the argument just given applies
with \(E\) and \(F\) exchanged, and in
the potentially good case the image of \(I_2\) in \(\GL_2(\Z_\ell)\) acting on \(T_\ell(F)\) is finite
\cite[Corollary~VII.7.3]{Si1}, nontrivial (I7) and contained in \(\SL_2(\Z_\ell)\); it has order prime to \(\ell\),
because an element of order \(\ell\) of \(\SL_2(\Q_\ell)\) would have a primitive \(\ell\)-th root of unity as an
eigenvalue, of degree \(\ell-1>2\) over \(\Q_\ell\); and it injects into \(\GL_2(\F_\ell)\), as the kernel of
reduction of \(\GL_2(\Z_\ell)\) is torsion-free. So \(I_2\) acts on \(F[\ell]\) through a nontrivial group of order
prime to \(\ell\), while it acts on \(E[\ell]\) through a group of order \(1\) or \(\ell\). If \(F\) had good reduction at \(2\),
then by \hyperref[lemma-localtorsion]{Lemma~7}(c) \(4\mid\#\widetilde F(\F_2)\le5\), so \(a_2(F)=-1\), while
\(t_2(E)=\pm3\), and \(\ell\nmid-1\mp3\). So \(F\) is multiplicative at \(2\). A nonsplit \(F\) at \(2\) has at most
\(4\) rational \(2\)-primary points (\hyperref[lemma-localtorsion]{Lemma~7}(b)), so \(F\) is split at \(2\), and
then \(E\) is split at \(2\) since \(3-(-3)=6\) is prime to \(\ell\). A nonsplit \(E\) at \(3\) has no rational point
of order \(3\) (Lemma~7(b): the torsion order divides \(8\)), so \(E\) is split at \(3\), and then \(F\) is split at
\(3\) since \(4-(-4)=8\) is prime to \(\ell\).

\emph{The primes \(5\), \(7\), \(11\), \(13\).} The prime-to-\(p\) torsion of \(E\) and \(F\) has order \(12\) and
\(8\), so neither curve is additive (\hyperref[lemma-localtorsion]{Lemma~7}(a)). If \(X\in\{E,F\}\) has good
reduction, \(T_X\) injects into \(\widetilde X(\F_p)\) (Lemma~7(c)), so \(12\), resp.~\(8\), divides
\(p+1-a_p(X)\); with \(a_p^2\le4p\) this leaves the traces
\[\begin{array}{c|cccc} p&5&7&11&13\\ \hline a_p(E)&-&-4&0&2\\ a_p(F)&-2&0&\pm4&-2,\ 6\end{array}\]
(certificate, Part~C, which also checks that these are exactly the traces of the curves over \(\F_p\) with the
corresponding subgroup). A nonsplit \(E\) at \(7\) or \(13\) has no point of order \(3\) (Lemma~7(b), as
\(3\nmid16,28\)), so \(t_p(E)=p+1\) there. For \(\ell\neq p\): every difference between a good trace of one curve and
any possible value of \(t_p\) of the other is a nonzero integer whose prime factors are \(2\) and \(3\); the sets of
differences are \(\{4,8\}\), \(\{4,8,12\}\), \(\{4,8,12,16\}\), \(\{4,8,12,16\}\) for \(p=5,7,11,13\) (certificate,
Part~C). So by \hyperref[lemma-compare]{Lemma~8}(a) both curves are multiplicative at \(p\). For \(\ell=p\) we use
Lemma~8(b): good reduction of \(E\) is supersingular exactly at \(11\) and good reduction of \(F\) exactly at \(7\);
the values \(u\) of ordinary good reductions are \(3\) for \(F\) at \(5\), \(3\) for \(E\) at \(7\), \(\pm4\) for
\(F\) at \(11\), and \(2\) for \(E\), \(6\) and \(11\) for \(F\) at \(13\), none of them \(\pm1\), and those of \(E\)
and \(F\) at the same prime are different. So in every case both curves are multiplicative at \(p\).

Splitting types: at \(5\) and \(13\), a nonsplit \(F\) would have at most \(4\) rational \(2\)-primary points
(Lemma~7(b): \(2(p+1)=12,28\)), so \(F\) is split; at \(7\) and \(13\), \(E\) is split by its point of order \(3\). At
\(13\) both are therefore split. At \(5\) (for \(E\)), \(7\) (for \(F\)) and \(11\) (for both), opposite splitting
types would give \(t_p(E)-t_p(F)=\pm2(p+1)\in\{\pm12,\pm16,\pm24\}\), prime to \(\ell\) if \(\ell\neq p\), and
different unramified constituents \(\pm1\) if \(\ell=p\). The conductor statement follows from
\hyperref[lemma-gluing]{Lemma~1}(d) and (I13). \qed

\section{The group \([2,28]\)}\label{sec-228}

\begin{statement}{Theorem 11 (the group [2,28]).}\label{thm-228}
Let \(A\), \(E\), \(F\), \(n\) be as in Section~\ref{gluing} with \(A(\Q)_\tors\cong\Z/2\times\Z/28\) and \(n\ge2\).
Then \(E\) and \(F\) have split multiplicative reduction at \(3\) and at \(5\), so \(v_3(N_A)=v_5(N_A)=2\).
\end{statement}

\emph{Proof.} By \hyperref[prop-parents]{Proposition~3}(c), \(T_E=\Z/7\), \(T_F=\Z/2\times\Z/4\), and \(n\) is prime
to \(14\); let \(\ell\) be a prime factor of \(n\), so \(\ell\) is odd, \(\ell\neq7\), and \(E[\ell]\cong
F[\ell]\). By \hyperref[lemma-localtorsion]{Lemma~7}(a), neither curve is additive at \(3\) or \(5\) (\(7>4\) and
\(8>4\)). Good reduction forces \(7\mid\#\widetilde E(\F_p)\), resp.~\(8\mid\#\widetilde F(\F_p)\) (Lemma~7(c)); with
the Hasse bound this leaves \(a_3(E)=-3\), no good reduction of \(F\) at \(3\), \(a_5(E)=-1\) and \(a_5(F)=-2\)
(certificate, Part~D). A nonsplit \(E\) at \(3\) or \(5\) has no point of order \(7\) (Lemma~7(b):
\(7\nmid8,12\)), and a nonsplit \(F\) at \(5\) has at most \(4\) rational \(2\)-primary points
(\(2(5+1)=12\)).

\emph{At \(3\).} \(F\) is multiplicative. If \(E\) were good, then \(\ell\neq3\) would divide \(-3-4=-7\) or
\(-3+4=1\), impossible as \(\ell\neq7\); and \(\ell=3\) is impossible because \(E\) is then supersingular
(\(a_3=-3\equiv0\)) and \(F\) is not (\hyperref[lemma-compare]{Lemma~8}). So \(E\) is multiplicative, hence split. If
\(F\) were nonsplit, \(t_3(E)-t_3(F)=8\) would be divisible by \(\ell\neq3\), or the unramified constituents \(1\) and
\(-1\) would differ if \(\ell=3\). So \(F\) is split.

\emph{At \(5\).} If \(E\) were good, then for \(\ell\neq5\) the differences \(-1-(-2)=1\) and \(-1-6=-7\) leave no
\(\ell\), and for \(\ell=5\) the unramified constituents \(4\) of \(E\) and \(3\) (good \(F\)) or \(1\) (split \(F\))
differ. So \(E\) is multiplicative, hence split. If \(F\) were good, then \(6-(-2)=8\) leaves no \(\ell\neq5\), and
the constituents \(1\) and \(3\) differ for \(\ell=5\); so \(F\) is multiplicative, hence split. The conductor follows
from \hyperref[lemma-gluing]{Lemma~1}(d). \qed

\section{The group \([32]\)}\label{sec-32}

\begin{statement}{Theorem 12 (two curves with torsion Z/8).}\label{thm-88}
Let \(A\), \(E\), \(F\), \(n\) be as in Section~\ref{gluing} with \(A(\Q)_\tors\cong\Z/32\) and \(T_E\cong
T_F\cong\Z/8\). Then \(E\) and \(F\) are not isogenous over \(\bar\Q\).
\end{statement}

\emph{Proof.} By \hyperref[prop-parents]{Proposition~3}(b), \(16\mid n\), so \(E[16]\cong F[16]\). If \(E\) and \(F\)
are isogenous over \(\Q\), then \(A\) is isogenous over \(\Q\) to \(E^2\), contradicting
\hyperref[thm-isotypic]{Theorem~5}(b). Suppose they are isogenous over \(\bar\Q\) but not over \(\Q\). By (I2) they
have no complex multiplication, so \(\Hom_{\bar\Q}(E,F)=\Z\varphi\), and \(g\varphi g^{-1}=\chi(g)\varphi\) for a
nontrivial quadratic character \(\chi\); hence \(V_2(F)\cong V_2(E)\otimes\chi\) and \(\tr\rho_F(g)=\chi(g)\tr
\rho_E(g)\) on \(V_2\). In a basis of \(T_2(E)\) whose first vector maps to a rational point of order \(8\), every
\(g\) acts by \(\left(\begin{smallmatrix}1&*\\0&\chi_2(g)\end{smallmatrix}\right)\) modulo \(8\), so \(\tr\rho_E(g)
\equiv1+\chi_2(g)\pmod 8\), and likewise for \(F\). For \(\chi(g)=-1\) this gives \(2(1+\chi_2(g))\equiv0\pmod8\),
that is \(\chi_2(g)\equiv-1\pmod4\) (certificate, Part~E). So the complement of \(\ker\chi\) is contained in that of
the kernel of the character \(\chi_{-4}\) of \(\Q(i)\); both kernels have index \(2\), so \(\chi=\chi_{-4}\). By
Dirichlet's theorem there is a prime \(p\equiv3\pmod8\) of good reduction for \(E\) and \(F\). Then \(a_p(F)=-a_p(E)\), and \(8\mid
\#\widetilde E(\F_p)\) gives \(a_p(E)\equiv p+1\equiv4\pmod8\), so \(a_p(E)-a_p(F)=2a_p(E)\equiv8\pmod{16}\). This
contradicts \(a_p(E)\equiv a_p(F)\pmod{16}\) (\hyperref[lemma-gluing]{Lemma~1}(c)). \qed

\begin{statement}{Theorem 13 (torsion Z/8 and Z/4 with the same j-invariant).}\label{thm-samej}
Let \(A\), \(E\), \(F\), \(n\) be as in Section~\ref{gluing} with \(A(\Q)_\tors\cong\Z/32\), \(T_E\cong\Z/8\) and
\(T_F\cong\Z/4\). Then \(j(E)\neq j(F)\).
\end{statement}

\emph{Proof.} Suppose \(j(E)=j(F)\). By (I2), \(E\) has no complex multiplication; fix an isomorphism
\(\rho\colon E\to F\) over \(\bar\Q\). Then \(g\rho g^{-1}=\chi(g)\rho\) for a quadratic character \(\chi\), with
field \(K\); \(\chi\neq1\) because \(T_E\not\cong T_F\). Let \(P_4\in T_E\) have order \(4\), let \(Q'\) generate
\(T_F\), and \(Q_4=\rho^{-1}(Q')\), so that \(gQ_4=\chi(g)Q_4\). Then \(2Q_4\) is a rational point of order \(2\),
hence the unique one, \(2P_4\), and \(Q_4-P_4\in E[2]\). For \(\chi(g)=-1\), \(g(Q_4-P_4)=(Q_4-P_4)+2P_4\), and for
\(\chi(g)=1\) it is fixed; so \(Q_4-P_4\) is a point of order \(2\) with field of definition \(K\), and
\(K=\Q(E[2])\). Thus \(E(K)\supseteq\Z/2\times\Z/8\), and by (I3) \(E(K)_\tors\cong\Z/2\times\Z/8\) has exponent
\(8\).

By \hyperref[prop-parents]{Proposition~3}(b), \(n=4m\) with \(m\) odd. Let \(\psi_4\colon E[4]\to F[4]\) be the
restriction of \(\psi\), an anti-isometry, and \(S=\rho^{-1}\psi_4\), so that \(\det S\equiv-1\pmod4\) and
\(Sg=\chi(g)gS\) on \(E[4]\). In a basis of \(E[4]\) beginning with \(P_4\), complex conjugation is
\(\left(\begin{smallmatrix}1&b\\0&-1\end{smallmatrix}\right)\) with \(b\in\{0,1\}\) (as in the proof of
\hyperref[thm-isotypic]{Theorem~5}), and \(\chi(c)=(-1)^b\) since \(K=\Q(E[2])\). The solutions of
\(Sc=\chi(c)cS\), \(\det S=-1\) in \(M_2(\Z/4)\) are \(8\) for \(b=0\) and \(4\) for \(b=1\) (of the form
\(\left(\begin{smallmatrix}x&y\\2x&-x\end{smallmatrix}\right)\) with \(x\) odd and \(y\) even), and all are \(\equiv I
\pmod2\) (certificate, Part~E). Hence \(\psi\) and \(\rho\) agree on \(E[2]\).

The quotient \(A_4=(E\times F)/\Gamma_{\psi_4}\) maps onto \(A\) with kernel \(\Gamma_\psi/\Gamma_{\psi_4}\) of odd
order \(m^2\), so \(A_4(\Q)[2^\infty]\cong A(\Q)[2^\infty]\) contains a point of order \(32\). Let
\(B=(E\times F)/\Gamma_{\rho|E[2]}\). Over \(K\), \((P,Q)\mapsto(P+\rho^{-1}Q,\ P-\rho^{-1}Q)\) is an isogeny
\(E\times F\to E\times E\) with kernel \(\Gamma_{\rho|E[2]}\) (its matrix \(\left(\begin{smallmatrix}1&1\\1&-1
\end{smallmatrix}\right)\) has determinant \(-2\)), so it induces an isomorphism \(B_K\cong E_K\times E_K\); an
element \(g\) with \(\chi(g)=-1\) changes \(\rho^{-1}\) into \(-\rho^{-1}\), that is, it exchanges the two factors.
Hence, for a fixed \(g\) with \(\chi(g)=-1\), \(B(\Q)\cong\{(x,gx):x\in E(K)\}\cong E(K)\), whose torsion has
exponent \(8\). Since \(\psi\) and \(\rho\)
agree on \(E[2]\), \(\Gamma_{\rho|E[2]}\subseteq\Gamma_{\psi_4}\), and the induced isogeny \(\beta\colon B\to A_4\) has
kernel \(\Gamma_{\psi_4}/\Gamma_{\rho|E[2]}\), killed by \(2\); so there is \(\beta'\colon A_4\to B\) over \(\Q\) with
\(\beta'\beta=[2]\) and \(\beta\beta'=[2]\), and \(\ker\beta'\subseteq A_4[2]\). A rational point \(R\) of order
\(32\) on \(A_4\) gives \(\beta'(R)\in B(\Q)\) of order at least \(16\): a contradiction. \qed

\medskip
Together with \hyperref[prop-parents]{Proposition~3}(b) and \hyperref[thm-isotypic]{Theorem~5}(b) this proves
part~(d) of the Theorem.

\begin{statement}{Remark 1 (what remains for [32]).}\label{rem-32}
For a principally polarized abelian surface over \(\Q\), split over \(\Q\), with rational torsion \(\Z/32\), the
complementary curves fall into one of three cases, none of which is excluded here:
\begin{itemize}
\item \textbf{(A)} \(T_E\cong\Z/8\), \(T_F\cong\Z/4\), \(v_2(n)=2\), and \(E\), \(F\) are not isogenous over \(\bar\Q\);
\item \textbf{(B)} \(T_E\cong\Z/8\), \(T_F\cong\Z/4\), \(v_2(n)=2\), and \(F\) is isogenous over \(\Q\) to a nontrivial
quadratic twist of \(E\), with \(j(F)\neq j(E)\);
\item \textbf{(C)} \(T_E\cong T_F\cong\Z/8\), \(16\mid n\), and \(E\), \(F\) are not isogenous over \(\bar\Q\).
\end{itemize}
In case (B), \(E\) has no complex multiplication, so a curve isogenous to \(E\) over \(\bar\Q\) is isogenous over
\(\Q\) to a quadratic twist of \(E\), and the trivial twist is excluded by \hyperref[thm-isotypic]{Theorem~5}.
\end{statement}

\section{The group \([25]\)}\label{sec-25}

In this section \(A\), \(E\), \(F\), \(n\) are as in Section~\ref{gluing} with \(A(\Q)_\tors\cong\Z/25\); by
\hyperref[prop-parents]{Proposition~3}(a), \(T_E\cong T_F\cong\Z/5\) and \(5\mid n\). For \(d\in\Q^\times\) with
\(d^2-11d-1\neq0\) let \(E_d\colon y^2+(1-d)xy-dy=x^3-dx^2\) (I4); its invariants are \(\Delta=d^5(d^2-11d-1)\),
\(c_4=d^4-12d^3+14d^2+12d+1\) and \(c_6=-d^6+18d^5-75d^4-75d^2-18d-1\), and
\((x,y)\mapsto(x/d^2-1/d,\ y/d^3-x/d^2)\) is an isomorphism \(E_d\to E_{-1/d}\) for every \(d\neq0\) (certificate,
Part~F).

\begin{statement}{Lemma 14 (local behaviour of a rational point of order 5).}\label{lemma-five}
Let \(X\) be an elliptic curve over \(\Q\) with a rational point \(P\) of order \(5\).
\begin{enumerate}
\item[(a)] Let \(p\neq5\). Then \(X\) is not additive at \(p\); if it is nonsplit multiplicative at \(p\), then
\(5\mid p+1\); if it is good at \(p\), then \(5\mid\#\widetilde X(\F_p)\).
\item[(b)] Write \((X,P)\cong(E_d,(0,0))\). If \(v_5(d)\neq0\), \(X\) is split multiplicative at \(5\). If
\(v_5(d)=0\) and \(d\not\equiv3\pmod5\), \(X\) is good at \(5\), with \(\#\widetilde X(\F_5)=10\) if \(d\equiv2\)
and \(5\) if \(d\equiv1,4\pmod5\). If \(d\equiv3\pmod5\), \(X\) has additive, potentially good reduction at \(5\);
the minimal discriminant has valuation \(2\) if \(d\not\equiv18\pmod{25}\), and valuation \(3\), with
\(v_5(c_6)\ge2\), if \(d\equiv18\pmod{25}\).
\item[(c)] Suppose \(X\) is good or split multiplicative at \(5\). Then \(X[5]=C\oplus\langle P\rangle\) with \(C\)
the unique \(G_{\Q_5}\)-stable line on which \(G_{\Q_5}\) acts by \(\bar\chi_5\). Identify \(X[5]/C\) with \(\Z/5\)
through \(P\mapsto1\), and let \(h_X\in H^1(\Q_5,X[5]/C)=\Hom(G_{\Q_5},\Z/5)\) be the image of the Kummer class
\(\delta(P)\in H^1(\Q_5,X[5])\). If \(X\) is split multiplicative, \(h_X=0\). If \(X\) is good, \(h_X\) is
unramified and takes the value \(c(X)=(\alpha-1)/5\bmod5\) at Frobenius, where \(\alpha\in\Z_5^\times\) is the
unit root of \(Y^2-a_5(X)Y+5\); \(c(X)=4\) if \(\#\widetilde X(\F_5)=5\), and \(c(X)=3\) if
\(\#\widetilde X(\F_5)=10\).
\end{enumerate}
\end{statement}

\emph{Proof.} (a) \hyperref[lemma-localtorsion]{Lemma~7}.

(b) If \(v_5(d)>0\), the equation is integral, \(c_4\equiv1\) and \(v_5(\Delta)=5v_5(d)>0\), so the reduction is
multiplicative (I6); the reduced curve is \(y^2+xy=x^3\), with tangent lines \(y=0\) and \(y=-x\) at the node, so it
is split. If \(v_5(d)<0\), use the isomorphism \(E_d\cong E_{-1/d}\). If \(v_5(d)=0\), then \(d^2-11d-1=(d-3)^2-5(d+2)\), so
\(\Delta\) is a unit unless \(d\equiv3\pmod5\), and the counts are computed over \(\F_5\) (certificate, Part~F). If
\(d=3+5k\), then \(d^2-11d-1=25(k^2-k-1)\) and \(k^2-k-1=(k-3)^2+5(k-2)\), so \(v_5(\Delta)=2\) unless
\(k\equiv3\pmod5\); for \(d=18+25k'\), \(d^2-11d-1=125(1+5k'+5k'^2)\) and \(v_5(\Delta)=3\). Moreover
\(5\mid c_4\), \(5\mid c_6\), and \(25\mid c_6\) if \(d\equiv18\pmod{25}\) (certificate, Part~F). Since
\(v_5(\Delta)<12\) the equation is minimal; it is additive since \(5\mid c_4\), and \(v_5(j)=3v_5(c_4)-v_5(\Delta)\ge0\)
gives potentially good reduction (I6).

(c) In the split case, \(X\cong E_q\) with \(q=r^5\) for some \(r\in\Q_5^\times\): if \(P\) is represented by
\(w\in\Q_5^\times\), then \(w^5=q^k\) with \(5\nmid k\) because \(\mu_5(\Q_5)=1\), and \(r=w^aq^b\) for \(ak+5b=1\).
So \(X[5]=\mu_5\oplus\langle P\rangle\), \(C=\mu_5\), and the \(5\)-divisions \(U\) of \(P\) are represented by
fifth roots of \(w\), so \(gU-U\in\mu_5=C\) and \(h_X=0\). In the good case, \(P\) reduces to a point of order \(5\)
(\hyperref[lemma-localtorsion]{Lemma~7}(c)), so \(\widetilde X\) is ordinary \cite[Theorem~V.3.1]{Si1}, and by (I9)
the kernel \(C\) of \(X[5]\to\widetilde X[5](\bar\F_5)\cong\Z/5\) is a line with action \(\bar\chi_5\) (the
determinant), not containing \(P\). The lines \(C\) and \(\langle P\rangle\) are the only stable ones, since
\(\bar\chi_5\neq1\) on \(G_{\Q_5}\). For \(U\) with \(5U=P\), \(h_X(g)\) is \(g\widetilde U-\widetilde U\in
\widetilde X[5](\bar\F_5)=\langle\widetilde P\rangle\), which depends only on the image of \(g\) in \(G_{\F_5}\).
The group \(\widetilde X[5^\infty](\bar\F_5)\) is isomorphic to \(\Q_5/\Z_5\) \cite[Theorem~V.3.1(b)]{Si1}, and the
Frobenius endomorphism \(\pi\), which satisfies \(\pi^2-a_5\pi+5=0\) \cite[Theorem~V.2.3.1]{Si1}, acts on it by a
unit root of \(Y^2-a_5Y+5\), that is by \(\alpha\). With \(\widetilde P=1/5\) and \(\widetilde U=1/25+j/5\), and
\(\alpha\equiv a_5\equiv1\pmod5\), \(\pi\widetilde U-\widetilde U\equiv(\alpha-1)/25=c(X)\widetilde P\). The values
of \(c(X)\) are computed modulo \(25\) in Part~F of the certificate. \qed

\begin{statement}{Lemma 15 (the Kummer relation and the trace identity).}\label{lemma-kummer}
Let \(R\) generate \(A(\Q)_\tors\cong\Z/25\), let \(P=\pi_E(R)\) and \(Q=\pi_F(R)\).
\begin{enumerate}
\item[(a)] \(P\) and \(Q\) have order \(5\), and \(Q=u\,\psi(P)\) for a unit \(u\in(\Z/5)^\times\), with \(u=1\) if
and only if \(25\mid n\).
\item[(b)] \(\psi_*\delta_E(P)=u\,\delta_F(\psi(P))\) in \(H^1(\Q,F[5])\), where \(\delta_E\colon E(\Q)\to
H^1(\Q,E[5])\) and \(\delta_F\) are the Kummer maps.
\item[(c)] Put \(\lambda=u^{-1}\). For every \(g\in G_{\Q}\),
\[\tr\rho_F(g)-1-\chi_5(g)\equiv\lambda\bigl(\tr\rho_E(g)-1-\chi_5(g)\bigr)\pmod{25},\]
where \(\rho_E\), \(\rho_F\) are the representations on \(T_5(E)\), \(T_5(F)\). In particular
\(\#\widetilde F(\F_p)\equiv\lambda\,\#\widetilde E(\F_p)\) modulo \(25\) at every prime \(p\neq5\) of good reduction for
\(E\) and \(F\), and \(\lambda\in\{2,3,4\}\) if \(v_5(n)=1\).
\end{enumerate}
\end{statement}

\emph{Proof.} (a) If \(\pi_E(R)=0\), then \(R\in F(\Q)\), which has no point of order \(25\) by (I1); and
\(5\pi_E(R)=0\) as \(T_E=\Z/5\). So \(P\) has order \(5\), and so has \(Q\). As \(\psi\) is equivariant, \(\psi(P)\)
generates \(T_F\cong\Z/5\), so \(Q=u\psi(P)\) with \(u\neq0\). Since \(i_F(\psi(x))=-i_E(x)\) for \(x\in E[n]\)
(I10), \(nR=\sigma\alpha(R)=i_E(P)+u\,i_F(\psi(P))=(1-u)i_E(P)\). If \(25\mid n\), \(nR=0\) and \(u=1\); if
\(v_5(n)=1\), \(nR=(n/5)\cdot5R\neq0\), so \(u\neq1\).

(b) Write \(R=\sigma(U,V)\); by \hyperref[lemma-gluing]{Lemma~1}(e), \(nU=P\), \(nV=Q\) and \(gV-V=\psi(gU-U)\). With
\(m=n/5\), \(mU\) and \(mV\) are \(5\)-divisions of \(P\) and \(Q\), and \(g(mV)-mV=\psi(g(mU)-mU)\). So
\(\delta_F(Q)=\psi_*\delta_E(P)\), and \(\delta_F(Q)=u\,\delta_F(\psi(P))\).

(c) Choose a basis \((e_1,e_2)\) of \(T_5(E)\) with \(e_1/5=P\), and \((f_1,f_2)\) of \(T_5(F)\) with
\(f_i/5=\psi(e_i/5)\). Then \(\rho_E(g)=\left(\begin{smallmatrix}1+5a&b\\5c&d\end{smallmatrix}\right)\) and
\(\rho_F(g)=\left(\begin{smallmatrix}1+5a'&b'\\5c'&d'\end{smallmatrix}\right)\) with \(b\equiv b'\), \(d\equiv d'\pmod5\)
and \(\det\rho_E(g)=\det\rho_F(g)=\chi_5(g)\). For such matrices, exactly,
\[\tr M-1-\det M=5\bigl(a(1-d)+bc\bigr)\]
(certificate, Part~G), so \(\tr\rho_E(g)-1-\chi_5(g)\equiv5(a(1-d)+bc)\) and \(\tr\rho_F(g)-1-\chi_5(g)\equiv
5(a'(1-d)+bc')\pmod{25}\). The cocycle \(g\mapsto(\rho_E(g)-1)e_1/25=(ae_1+ce_2)/5\) represents \(\delta_E(P)\), and
likewise \((a',c')\) represents \(\delta_F(\psi(P))\) in the basis \(f_i/5\). By (b) there is \(w\in F[5]\), with
second coordinate \(w_2\), such that \(u(a',c')\equiv(a,c)+(\rho_F(g)-1)w\), that is \(ua'\equiv a+bw_2\) and
\(uc'\equiv c+(d-1)w_2\pmod5\). Then
\[u\bigl(a'(1-d)+bc'\bigr)\equiv(a+bw_2)(1-d)+b(c+(d-1)w_2)=a(1-d)+bc\pmod5\]
(certificate, Part~G), which is the identity. At a prime \(p\neq5\) of good reduction, a Frobenius element has trace
\(a_p\) and \(\chi_5=p\), so \(\tr\rho-1-\chi_5=-\#\widetilde X(\F_p)\). The last claim follows from (a). \qed

\medskip
The identity holds for all \(g\in G_{\Q}\), in particular for Frobenius lifts at primes of bad reduction: if
\(E\) is split multiplicative at \(p\neq5\), any Frobenius lift has trace \(p+1\) on \(T_5(E)\) (I8), so if \(F\) is
good at \(p\), then \(25\mid\#\widetilde F(\F_p)\); and symmetrically. \(\quad(\ast)\)

\begin{statement}{Lemma 16 (additive reduction at 5).}\label{lemma-additive}
\begin{enumerate}
\item[(a)] Let \(B/\Q_5\) be an abelian surface isogenous to \(E'\times F'\), where \(E'\) and \(F'\) are elliptic
curves with additive, potentially good reduction. Then \(B(\Q_5)\) has no point of order \(25\).
\item[(b)] Let \(d\in\Z_5\) with \(d\equiv3\pmod5\) and \(v_5(2d-11)\le2\), and let \(K=\Q_5(\zeta_5)\). Then
\(E_d[5]\not\subseteq E_d(K)\).
\item[(c)] Let \(d\in\Z_5\) with \(d\equiv18\pmod{25}\) (for instance \(d\equiv68\pmod{125}\)), and \(P=(0,0)\). Then
the Kummer class \(\delta(P)\in H^1(\Q_5,E_d[5])\) is represented by a homomorphism \(G_{\Q_5}\to\langle P\rangle\)
that is nonzero on \(I_5\).
\end{enumerate}
\end{statement}

\emph{Proof.} (a) Let \(M=\Q_5(\zeta_3)\), unramified over \(\Q_5\), and \(L=M(5^{1/12})\), cyclic of degree \(12\)
and totally ramified over \(M\); let \(\tau\) generate \(\Gal(L/M)\). An elliptic curve over \(\Q_5\) with
\(v_5(j)\ge0\) has good reduction over \(L\): if \(\varpi=5^{1/12}\) and \(\nu=v_5(\Delta)\) for a minimal model, the
equation \(y^2=x^3-27c_4\varpi^{-4\nu}x-54c_6\varpi^{-6\nu}\) is integral, because \(3v_5(c_4)\ge\nu\) and
\(2v_5(c_6)\ge\nu\) by \(c_6^2=c_4^3-1728\Delta\), and its discriminant \(6^{12}\Delta\varpi^{-12\nu}\) is a unit. So
\(E'\), \(F'\) and \(B\) have good reduction over \(L\) (I7). The group \(I_M\) acts on \(V_\ell(E')\)
(\(\ell\neq5\)) through \(\Gal(L/M)\), and nontrivially, since \(E'\) is additive over \(M\) (I6), (I7); as
\(\det\rho_\ell=\chi_\ell\) is trivial on \(I_M\), \(\rho_\ell(\tau)\) is a nontrivial element of finite order of
\(\SL_2(\Q_\ell)\), so it has no eigenvalue \(1\) and \(V_\ell(E')^{I_M}=0\). The same holds for \(F'\), so
\(\sum_{i<12}\rho_\ell(\tau)^i\), whose image consists of \(\tau\)-invariant vectors, is \(0\) on \(V_\ell(B)\cong
V_\ell(E')\oplus V_\ell(F')\). By the N\'eron mapping property, \(\Gal(L/M)\) acts on the N\'eron model of \(B\)
over the ring of integers of \(L\) (see \cite[\S2]{ST}); as \(L/M\) is totally ramified, \(\tau\) acts on the
reduction \(\widetilde B\) of \(B\) over \(L\) by an automorphism \([\tau]\), compatibly with reduction and with
\(V_\ell(B)\cong V_\ell(\widetilde B)\). So \(\sum_{i<12}[\tau]^i=0\), as \(V_\ell\) is faithful on endomorphisms. A point \(R\in B(\Q_5)\) of \(5\)-power order reduces to a point fixed by
\([\tau]\), hence killed by \(12\), hence zero. So \(R\) lies in the kernel of reduction, the group of points of a
formal group of dimension \(2\) over the ring of integers of \(L\) (I11), with \(e=v_L(5)=12<20\);
\hyperref[lemma-formal]{Lemma~6}(c) shows that \(R\) does not have order \(25\).

(b) Let \(\zeta\) be a primitive fifth root of unity, \(\sqrt5=\zeta-\zeta^2-\zeta^3+\zeta^4\),
\(\varepsilon=(\sqrt5-1)/2\), \(\alpha_0=(11+5\sqrt5)/2\) and \(\beta_0=(11-5\sqrt5)/2\), the roots of
\(d^2-11d-1\), and \(q(d)=(d-\alpha_0)/(d-\beta_0)\). Part~F of the certificate checks, in \(\Q(\zeta)(u)\) with
\(d(u)=-(\varepsilon^5u^5-\varepsilon^{-5})/(u^5+1)\): (i) \(u^5=-q(d(u))\); (ii) \(x(u)=N(u)/D(u)\), where
\(N(u)=-\varepsilon^4\prod_{j=0}^{3}(u-\zeta^j(1+\zeta)^2)\) and \(D(u)=(u+\zeta^2)(u+\zeta^3)(u+1)^2\), is a root of
the \(5\)-division polynomial of \(E_{d(u)}\); (iii) \(N\) and \(D\) are coprime of degree \(4\). These are the
formulas of Morton \cite{Mo}. For \(d=3+5k\) one has \(q(d)=(2k-1-\sqrt5)/(2k-1+\sqrt5)\). Normalize the valuation
\(v\) of \(K\) by \(v(5)=4\), so \(v(\sqrt5)=2\); note \(v_5(2d-11)=1+v_5(2k-1)\). If \(v_5(2k-1)=0\), then
\(q-1=-2\sqrt5/(2k-1+\sqrt5)\) has valuation \(2\). If \(2k-1=5w\) with \(w\) a unit, then
\(-q=(1-\sqrt5w)/(1+\sqrt5w)\) and \(v(-q-1)=2\). A unit \(y\equiv1\) that is a fifth power \(x^5\) in \(K\) has
\(x\equiv1\) (the residue field is \(\F_5\)), so \(v(y-1)\ge\min(4+v(x-1),5v(x-1))\ge5\). Hence \(-q(d)\) is not a
fifth power in \(K\) (in the first case neither is \(q(d)\), and \(-1\) is a fifth power). As \(\mu_5\subset K\), a
root \(u_0\) of \(u^5=-q(d)\) has degree \(5\) over \(K\); \(d(u_0)=d\), because \(d\mapsto-q(d)\) is a Möbius
transformation; \(D(u_0)\neq0\) since the roots of \(D\) lie in \(K\); and \(x(u_0)\) is the \(x\)-coordinate of a
point of \(E_d[5]\smallsetminus\{0\}\). If \(x(u_0)\) were in \(K\), \(u_0\) would be a root of the nonzero
polynomial \(N-x(u_0)D\), of degree at most \(4\) over \(K\). So \(x(u_0)\notin K\).

(c) By \hyperref[lemma-five]{Lemma~14}(b), \(E_d\) is additive with \(v_5(\Delta)=3\) and \(v_5(c_6)\ge2\). Let
\(L_4=\Q_5(5^{1/4})\), cyclic of degree \(4\) and totally ramified over \(\Q_5\) (as \(\mu_4\subset\Q_5\)), with
generator \(\tau\), and \(\varpi=5^{1/4}\). As in (a), the equation \(y^2=x^3-27(c_4/5)x-54c_6\varpi^{-6}\) is integral
with unit discriminant, and \(v(c_6\varpi^{-6})>0\), so \(E_d\) has good reduction over \(L_4\), with reduction
\(\widetilde E\colon y^2=x^3+ax\), \(a\neq0\), of \(j\)-invariant \(1728\); it is ordinary, as \(5\equiv1\pmod4\)
\cite[Example~V.4.5]{Si1} (certificate, Part~F). As in (a), \(\sum_{i<4}[\tau]^i=0\) on \(\widetilde E\), so \(P\)
reduces to a point killed by \(4\), that is to \(0\). Let \(E_1\subset E_d(\bar\Q_5)\) be the set of points that
reduce to \(0\) in this model; as \(x\) in the model above is \(x\varpi^{-2}\) for the coordinate \(x\) of
\(y^2=x^3-27c_4x-54c_6\), \(E_1\) consists of \(0\) and the points with \(v_5(x)<1/2\) there, so it is a
\(G_{\Q_5}\)-stable subgroup. By \cite[\S1.11]{Se} (good reduction of height \(1\), any ramification), the kernel of
reduction \(E_1\cap E_d[5]\) is cyclic of order \(5\) and \(E_d[5]\to\widetilde E[5](\bar\F_5)\) is onto; so
\(E_1\cap E_d[5]=\langle P\rangle\). Some \(5\)-division \(U\) of \(P\) lies in \(E_1\): if \(5U_0=P\), then
\(5\widetilde U_0=\widetilde P=0\), so \(\widetilde U_0\) is the reduction of some \(T\in E_d[5]\), and \(U=U_0-T\)
reduces to \(0\). For \(g\in G_{\Q_5}\), \(gU-U\in E_1\cap
E_d[5]=\langle P\rangle\), so \(h(g)=gU-U\) is a homomorphism \(G_{\Q_5}\to\langle P\rangle\) representing
\(\delta(P)\). If \(h\) vanished on \(I_5\), it would vanish on \(G_{M'}\) for the unramified extension \(M'\) of
\(\Q_5\) of degree \(5\), and \(U\) would be a point of order \(25\) in \(E_d(M')\). Over \(M'L_4\), of ramification
index \(4\), the argument above (with \(\Gal(M'L_4/M')\cong\Gal(L_4/\Q_5)\)) shows that \(U\) reduces to \(0\), so
\(U\) lies in a formal group over the ring of integers of \(M'L_4\), with \(e=4<20\), contradicting
\hyperref[lemma-formal]{Lemma~6}(c). \qed

\begin{statement}{Theorem 17 (the group [25]: reduction types).}\label{thm-25}
In the setting of this section:
\begin{enumerate}
\item[(i)] \(E\) and \(F\) are semistable;
\item[(ii)] at \(5\), either both have good reduction and \(\{\#\widetilde E(\F_5),\#\widetilde F(\F_5)\}=\{5,10\}\),
in which case \(v_5(n)=1\), or both have split multiplicative reduction; in particular the second case holds if
\(25\mid n\);
\item[(iii)] both \(E\) and \(F\) have split multiplicative reduction at \(2\) and \(3\), and either both have split
multiplicative reduction at \(11\), or both have split multiplicative reduction at \(5\) and \(7\). Hence \(66\)
divides both \(N_E\) and \(N_F\), or \(210\) divides both.
\end{enumerate}
\end{statement}

\emph{Proof.} Let \(R\), \(P\), \(Q\), \(u\) and \(\lambda\) be as in \hyperref[lemma-kummer]{Lemma~15}, so \(T_E=
\langle P\rangle\) and \(T_F=\langle\psi(P)\rangle\). Recall \(E[5]\cong F[5]\) (\(5\mid n\)). We use repeatedly:
if \(E\) and \(F\) have good reduction at a prime \(p\) and \(\#\widetilde E(\F_p)=\#\widetilde F(\F_p)=N\) with
\(v_5(N)=1\), we get a contradiction. Indeed both Frobenius polynomials equal \(f(Y)=Y^2-(p+1-N)Y+p\), so
\(f(\varphi)=0\) on \(V_\ell(E)\) and \(V_\ell(F)\) (\(\ell\neq p\), \(\varphi\) a Frobenius lift), hence on
\(V_\ell(A)\) (\hyperref[lemma-gluing]{Lemma~1}(d)); \(A\) has good reduction at \(p\) (I7), \(\varphi\) acts on
\(V_\ell(A)\cong V_\ell(\widetilde A)\) as the Frobenius endomorphism \(\pi\), so \(f(\pi)=0\) and \(f(1)=N\) kills
\(\widetilde A(\F_p)\). But \(A(\Q)[5^\infty]\) injects into \(\widetilde A(\F_p)\): by (I11) if \(p\neq5\), and by
(I11) and \hyperref[lemma-formal]{Lemma~6}(c) (\(e=1<4\)) if \(p=5\). So \(A(\Q)[5^\infty]\) would be killed by \(5\).

\emph{Primes \(p\neq5\).} By \hyperref[lemma-five]{Lemma~14}(a), neither curve is additive.

\emph{The prime \(5\).} By Lemma~14(b), each of \(E\), \(F\) is good, split multiplicative, or additive with a
parameter \(d\equiv3\pmod5\). In the good and split cases let \(C_E\subset E[5]\), \(C_F\subset F[5]\) be the lines of
Lemma~14(c). We push the relation \(\psi_*\delta_E(P)=u\,\delta_F(\psi(P))\) of Lemma~15(b) to
\(H^1(\Q_5,F[5]/C_F)=\Hom(G_{\Q_5},\Z/5)\), identifying \(F[5]/C_F\) with \(\Z/5\) through \(\psi(P)\mapsto1\).
\begin{itemize}
\item Both additive: \(A\) is isogenous over \(\Q_5\) to \(E\times F\), and \(R\) contradicts
\hyperref[lemma-additive]{Lemma~16}(a).
\item One additive, say \(E\cong E_d\), the other good or split (the other order is symmetric, as \(u\) is a unit):
by Lemma~14(c), \(F[5]\cong E[5]\) is the direct sum of \(\bar\chi_5\) and the trivial character, so \(G_K\) acts
trivially on \(E[5]\), \(K=\Q_5(\zeta_5)\). If \(v_5(2d-11)\le2\), this contradicts Lemma~16(b). Otherwise
\(d\equiv68\pmod{125}\); by Lemma~16(c), \(\delta_E(P)\) is represented by a homomorphism into \(\langle P\rangle\),
which \(\psi\) maps isomorphically onto \(F[5]/C_F\) (as \(\psi(P)\notin C_F\)); so the image of the left side is
nonzero on \(I_5\), while the image \(u\,h_F\) of the right side is unramified or zero by Lemma~14(c).
\item One good and one split, say \(E\) good: \(\psi(C_E)=C_F\) (the unique lines with action \(\bar\chi_5\)), so the
image of the left side is \(h_E\), which is nonzero, and that of the right side is \(u\,h_F=0\).
\item Both good: the counts lie in \(\{5,10\}\) (Lemma~14(b)) and are different by the remark at the beginning of
the proof. Evaluating the images at a Frobenius lift gives \(c(E)=u\,c(F)\), so \(\lambda=u^{-1}=c(F)/c(E)\) is
\(3\cdot4^{-1}=2\) or \(4\cdot3^{-1}=3\) in \(\F_5\) (certificate, Part~F); hence \(\lambda\neq1\) and \(v_5(n)=1\)
by Lemma~15(a).
\end{itemize}
This proves (ii), and with the primes \(p\neq5\), (i).

\emph{Small primes.} Let \(p\in\{2,3,7,11,13\}\). Bad reduction at \(p\) is split, since \(5\nmid p+1\)
(Lemma~14(a)). A good reduction has \(\#\widetilde X(\F_p)\in\{5\}\), \(\{5\}\), \(\{5,10\}\), \(\{10,15\}\),
\(\{10,15,20\}\) for \(p=2,3,7,11,13\), by Lemma~14(a) and the Hasse bound, so it is never divisible by \(25\)
(certificate, Part~G). By \((\ast)\) after Lemma~15, \(E\) and \(F\) are both good or both split at \(p\). At \(2\)
and \(3\) two good reductions would have the same count \(5\), which is excluded; so both curves are split at \(2\)
and \(3\). At \(11\), if both are good, their counts are \(10\) and \(15\) (equal counts are excluded), and
Lemma~15(c) gives \(\lambda=4\) (certificate, Part~G). Then at \(7\) two good reductions would have counts \(5\) and
\(10\) and give \(\lambda\in\{2,3\}\), so both curves are split at \(7\); and at \(5\) two good reductions would give
\(\lambda\in\{2,3\}\), so by (ii) both are split at \(5\). The conductors follow from (I13). \qed

\begin{statement}{Remark 2 (a second proof at 2 and 3).}\label{rem-route}
The mixed case at \(p\in\{2,3\}\) can also be excluded with the Kummer relation instead of \((\ast)\). Suppose \(E\)
is good at \(p\), so \(\#\widetilde E(\F_p)=5\), and \(F\) is split multiplicative at \(p\). Restrict
\hyperref[lemma-kummer]{Lemma~15}(b) to \(G_{\Q_p}\). The class \(\delta_E(P)\) is unramified, since the
\(5\)-division points of \(P\) generate an extension unramified at \(p\) \cite[Proposition~VIII.1.5(b)]{Si1}, and
nonzero, because \(E_1(\Q_p)\) is uniquely \(5\)-divisible and \(\widetilde P\notin5\widetilde E(\F_p)=0\); so
the left side is a nonzero unramified class (\(F[5]\cong E[5]\) is unramified). On the other side \(F\cong E_q\) with
\(q=r^5\), \(r\in\Q_p^\times\) (as \(\mu_5(\Q_p)=1\)), \(\psi(P)\) is represented by \(r^k\) with \(5\nmid k\), and
\(\delta_F(\psi(P))\) is the image of \(k\) times the Kummer class of \(r\) in \(H^1(\Q_p,\mu_5)=
\Q_p^\times/\Q_p^{\times5}\). It is zero if \(5\mid v_p(r)\), since every unit of \(\Q_2\) or \(\Q_3\) is a fifth
power, and its restriction to \(I_p\) is nonzero otherwise. Neither is a nonzero unramified class. The two-good case
is excluded as in the proof of \hyperref[thm-25]{Theorem~17}. The two arguments are independent of each
other.
\end{statement}

\begin{statement}{Lemma 18 (nonsplit multiplicative primes).}\label{lemma-nonsplit}
In the setting of this section, there is no prime \(p\) at which \(E\) has nonsplit multiplicative reduction with
\(5\nmid v_p(\Delta_{\min}(E))\) and \(v_5(p+1)=1\). The same holds for \(F\).
\end{statement}

\emph{Proof.} Suppose \(p\) is such a prime; \(p\neq5\) by \hyperref[thm-25]{Theorem~17}(ii). By (I8),
\(T_5(E)\cong\eta\otimes T_5(E_q)\) with \(v_p(q)=v_p(\Delta_{\min}(E))\) prime to \(5\), so some element of \(I_p\)
acts on \(E[5]\) by a transvection, \(E[5]^{I_p}\) is a line, and a Frobenius lift \(\varphi\) acts on it by
\(\eta(\varphi)\bar\chi_5(\varphi)=-p\equiv1\pmod5\). As \(F[5]\cong E[5]\), \(F[5]\) is ramified at \(p\), so \(F\) is
not good at \(p\) (I7), nor additive (\hyperref[lemma-five]{Lemma~14}(a)). If \(F\) were split, \(F[5]^{I_p}\) would
be the line \(\mu_5\) (I8), on which \(\varphi\) acts by \(p\equiv-1\pmod5\); but an isomorphism \(E[5]\cong F[5]\)
maps \(E[5]^{I_p}\) onto \(F[5]^{I_p}\) compatibly with \(\varphi\). So \(F\) is nonsplit, and on \(V_5(E)\) and
\(V_5(F)\) the characteristic polynomial of \(\varphi\) is \(f(Y)=(Y+1)(Y+p)\) (I8). By Cayley--Hamilton \(f(\varphi)\)
is \(0\) on \(V_5(E)\oplus V_5(F)\cong V_5(A)\), hence on \(A[5^\infty]=V_5(A)/T_5(A)\). Every point of
\(A(\Q)[5^\infty]\) is fixed by \(\varphi\), so it is killed by \(f(1)=2(p+1)\), which has \(5\)-adic valuation
\(1\): a contradiction with \(R\). \qed

\medskip
The condition concerns the curve \(E\) itself, not its isogeny class: an isogeny of degree \(5\) can multiply
\(v_p(\Delta_{\min})\) by \(5\).

\begin{statement}{Remark 3 (an infinite family).}\label{rem-family}
\hyperref[lemma-nonsplit]{Lemma~18} excludes curves that the conditions at \(2\), \(3\) and \(5\) of
\hyperref[thm-25]{Theorem~17} allow. For every integer \(t\), let \(d=2244+10830t\). Then \(E_d\) has a rational point
of order \(5\); it has split multiplicative reduction at \(2\) and \(3\), since \(6\mid d\), \(c_4\equiv1\) modulo
\(2\) and \(3\), and the reduction is \(y^2+xy=x^3\); it has good reduction at \(5\) with \(5\) points, since
\(d\equiv4\pmod5\) (\hyperref[lemma-five]{Lemma~14}(b)); and it has nonsplit multiplicative reduction of type
\(I_1\) at \(19\), where \(v_5(19+1)=1\): \(d\equiv78\pmod{361}\), so \(d\equiv2\pmod{19}\),
\(d^2-11d-1\equiv171=9\cdot19\pmod{361}\), \(19\nmid c_4(2)\), and \(-c_6(2)\) is not a square modulo \(19\). By
\cite[Lemma~V.5.2 and Theorem~V.5.3]{Si2} the reduction is split if and only if \(\gamma=-c_4/c_6\) is a square in
\(\Q_{19}\), and \(\gamma\equiv-c_6\) modulo squares because \(c_6^2\equiv c_4^3\pmod{19}\); so it is nonsplit. So no \(E_d\) is \(E\) or \(F\) for a surface as in this section (certificate, Part~H, which
also checks Magma's local data for \(d=78\) and \(d=2244\)).
\end{statement}

\begin{statement}{Theorem 19 (the group [25]: conductors).}\label{thm-25cond}
In the setting of this section, and relying on the completeness of Cremona's tables below conductor \(500000\)
(I14), \(N_E\ge8778\), \(N_F\ge8778\) and \(\max(N_E,N_F)\ge500000\). Consequently \(N_A=N_EN_F\ge4389000000\).
\end{statement}

\emph{Proof.} By \hyperref[thm-25]{Theorem~17}, \(N_E\) is squarefree and divisible by \(66\) or by \(210\); \(T_E
\cong\Z/5\); and by \hyperref[lemma-nonsplit]{Lemma~18}, \(E\) has no witness prime, that is, no prime \(p\) with
\(v_5(p+1)=1\) at which \(E\) is nonsplit multiplicative with \(5\nmid v_p(\Delta_{\min})\). Part~H of the certificate
runs through all curves in Cremona's tables with conductor divisible by \(6\), not divisible by \(25\), and below
\(500000\): there are \(1472509\) of them, in \(980367\) isogeny classes, and exactly \(464\) have torsion
\(\Z/5\); of these, \(155\) have conductor divisible by \(66\) or \(210\), and \(76\) of those have no witness prime,
the smallest conductor among them being \(8778\) (the curve 8778x1). Below \(8778\) the curves with torsion \(\Z/5\)
and conductor divisible by \(66\) or \(210\) are 1254k1, 1914o1 and 1914p1, with witness primes \(19\), \(29\),
\(29\), where they have type \(I_1\) and are nonsplit, and \(v_5(20)=v_5(30)=1\). So \(N_E\ge8778\), and likewise
\(N_F\ge8778\). The isogeny class of each of these three curves consists of two curves joined by an isogeny of
degree \(5\); the other curve (1254k2, 1914o2, 1914p2) has trivial rational torsion, and type \(I_5\) at the witness
prime, so it is excluded only because it has no rational point of order \(5\); no curve in these classes has a
rational point of order \(2\).

If \(N_E\) and \(N_F\) are both below \(500000\), both curves are among the \(464\). If they are not isogenous over
\(\Q\), Part~H shows that no \(\lambda\in\F_5^\times\) satisfies \(\#\widetilde F(\F_p)/5\equiv\lambda\,
\#\widetilde E(\F_p)/5\) modulo \(5\) at all primes \(7\le p\le59\) of good reduction for both: each of the \(107415\) pairs
of curves in different isogeny classes fails at some such prime. This contradicts
\hyperref[lemma-kummer]{Lemma~15}(c). If they are isogenous over \(\Q\) (for instance \(E\cong F\)), \(A\) is
isogenous over \(\Q\) to \(E^2\), which contradicts \hyperref[thm-isotypic]{Theorem~5}(a). So
\(\max(N_E,N_F)\ge500000\). Finally \(N_A=N_EN_F\) by \hyperref[lemma-gluing]{Lemma~1}(d), as \(E\) and \(F\) are
semistable, and \(N_A\ge8778\cdot500000\). \qed

\begin{statement}{Remark 4 (the trace test and the hypotheses of the Theorem for [25]).}\label{rem-trace}
Curves that are isogenous over \(\Q\) have the same point counts, so they pass the test of
\hyperref[lemma-kummer]{Lemma~15}(c) with \(\lambda=1\); this is why \hyperref[thm-isotypic]{Theorem~5}(a) is needed.
In the tables, the pair 165066d1, 165066d2 and the \(464\) pairs \((E,E)\) pass only with \(\lambda=1\). Parts
(i)--(iii) of \hyperref[thm-25]{Theorem~17}, Lemmas~15 and~18 and the bound \(\max(N_E,N_F)\ge500000\) use only a
rational point \(R\) of order \(25\); the torsion of \(E\) and \(F\) is then \(\Z/5\) or \(\Z/10\). The bound
\(N_E\ge8778\) uses \(T_E\cong\Z/5\): the curve 66c1 has torsion \(\Z/10\), conductor \(66\), split multiplicative
reduction at \(2\), \(3\) and \(11\), and no witness prime. Adding the \(32\) curves with torsion \(\Z/10\) of the
range of Part~H, none of the \(122743\) pairs in different isogeny classes passes the test (certificate, Part~H), so
\(\max(N_E,N_F)\ge500000\) holds as soon as \(A(\Q)\) has a point of order \(25\).
\end{statement}

\section{Further remarks}\label{remarks}

\begin{statement}{Remark 5 (products of elliptic curves).}\label{rem-products}
The hypothesis \(n\ge2\) of \hyperref[thm-22212]{Theorems~10} and~\hyperref[thm-228]{11} is necessary. For elliptic
curves \(E\), \(F\) over \(\Q\), the product \(E\times F\) with the product polarization is a principally polarized
abelian surface over \(\Q\), split over \(\Q\), with gluing degree \(1\) and torsion \(T_E\times T_F\). The curve
26b1 has torsion \(\Z/7\) and conductor \(26\), 15a1 has torsion \(\Z/2\times\Z/4\) and conductor \(15\), and 30a2
has torsion \(\Z/2\times\Z/6\) and conductor \(30\). So 26b1 \(\times\) 15a1 with the product polarization has
torsion \(\Z/2\times\Z/28\) but conductor \(390\), and 26b1 has good reduction at \(3\) and \(5\); and 30a2 \(\times\)
15a1 with the product polarization has torsion \((\Z/2)^3\times\Z/12\) and conductor \(450\), which is not divisible
by \(30030^2\) (certificate, Part~P). The proofs need a prime \(\ell\mid n\). For
\([2,2,20]\), 66c1 \(\times\) 15a1 has torsion \((\Z/2)^2\times\Z/20\), in accordance with
\hyperref[prop-parents]{Proposition~3}(d), which does not use \(n\ge2\). A product never has a rational point of order
\(25\) or \(32\) (certificate, Part~A).
\end{statement}

\begin{statement}{Remark 6 (where the hypotheses are used).}\label{rem-hyp}
The principal polarization, defined over \(\Q\), enters through (I10): it provides the complement \(F\) with
\(A/E\cong F\), the gluing \(\psi\), and the Kummer relation of \hyperref[lemma-gluing]{Lemma~1}(e). That \(\psi\) is
an anti-isometry is used only in \hyperref[thm-samej]{Theorem~13}. \hyperref[thm-isotypic]{Theorem~5} uses no
polarization. The condition \(n\ge2\) is used only in \hyperref[thm-22212]{Theorems~10} and~\hyperref[thm-228]{11} and in the
last sentence of \hyperref[prop-parents]{Proposition~3}.
The exact torsion group is used in \hyperref[prop-parents]{Proposition~3}, which determines \(T_E\) and \(T_F\), and
for \([25]\) as described in \hyperref[rem-trace]{Remark~4}. Part~(e)(iv) relies on the completeness of Cremona's
tables below conductor \(500000\) (I14); everything else is proved for all conductors.
\end{statement}

\section{Computations}\label{computations}

The certificate
\href{https://genus-2-torsion.github.io/data/knowledge/sources/qsplit_conditions_checks.m}{\texttt{qsplit\_conditions\_checks.m}}
(Magma V2.30-1, about \(30\) seconds) consists of assertions and ends with
\texttt{ALL\_QSPLIT\_CONDITIONS\_CHECKS\_PASSED}; a failed assertion stops it before that line. Magma's exit status is
\(0\) in both cases, so only the final line certifies a run. Its log is
\href{https://genus-2-torsion.github.io/data/knowledge/sources/qsplit_conditions_checks.log}{\texttt{qsplit\_conditions\_checks.log}}.
SHA-256 of the script:
\texttt{3fee4d4ed8ad5102\allowbreak 6d21db91e4ad840a\allowbreak 3f82a785d4c5bd9e\allowbreak a18283a40d31188d};
of the log:
\texttt{1044cef3965daad4\allowbreak e5237230414e90a4\allowbreak 6622943733e4c7de\allowbreak 16c207b34484793a}.

Part~A: \hyperref[prop-parents]{Proposition~3}: for each of the five groups, the pairs of groups from (I1) with the
properties of \hyperref[lemma-gluing]{Lemma~1}(a); the \(p\)-parts used for the gluing degrees; the elements of order
\(4\) of \(\Z/32\); products of groups from (I1). Part~B: \hyperref[lemma-hlp]{Lemma~9}: the families \(E_t\) and
\(F_t\), their torsion points, \(j\)-invariants and singular fibres; Kubert's families for \(\Z/2\times\Z/6\) and
\(\Z/2\times\Z/4\) and the parameter changes \(t=(s-3)/2\) and \(t=(4v+1)/(4v-1)\); the index \(24\); the identities
of the hand proof; the covering of \(\mathbb P^1(\Q_2)\) and \(\mathbb P^1(\Q_3)\) by discs. Part~C:
\hyperref[thm-22212]{Theorem~10}: the traces over \(\F_p\), \(p=5,7,11,13\), by Legendre sums and the halving
criterion and again from the group structure, the difference sets, the values \(u\), and the arithmetic of the
torsion bounds at \(2\), \(3\), \(5\), \(7\), \(11\), \(13\). Part~D: \hyperref[thm-228]{Theorem~11}: point counts of
all curves over \(\F_3\) and \(\F_5\), and the comparisons. Part~E: \hyperref[thm-isotypic]{Theorem~5}(b) and
\hyperref[thm-88]{Theorems~12} and~\hyperref[thm-samej]{13}: the fixed vectors of \(c\), the basis change, the corner
identities, the solutions of the matrix equation, the product model, the list (I3), the congruences. Part~F:
\hyperref[lemma-five]{Lemmas~14} and~\hyperref[lemma-additive]{16}: the invariants of \(E_d\), \(E_d\cong E_{-1/d}\),
the polynomial identities at \(5\), the counts over \(\F_5\), the ordinarity of \(j=1728\), the unit roots, and
Morton's formulas. Part~G: \hyperref[lemma-kummer]{Lemma~15} and \hyperref[thm-25]{Theorem~17}: the two polynomial
identities, the possible point counts, and the scalars \(\lambda\). Part~H: \hyperref[lemma-nonsplit]{Lemma~18},
\hyperref[thm-25cond]{Theorem~19} and \hyperref[rem-family]{Remarks~3} and~\hyperref[rem-trace]{4}: the family of
Remark~3, the scan of Cremona's tables, the witness primes, the isogeny classes of 1254k1, 1914o1, 1914p1, and the
test on the \(107415\) and \(122743\) pairs. Part~P: \hyperref[rem-products]{Remark~5}. Part~V (validation, not used
in the proofs): the trace relations used in \hyperref[thm-88]{Theorem~12}, for the curve \(E(b,b/u)\) with
\(b=(2u-1)(u-1)\) and \(u=13/18\) (Kubert's family for \(\Z/8\) \cite[Table~3]{Ku}; torsion \(\Z/8\)) and its twist by
\(-1\) (torsion \(\Z/4\)), at the \(76\) primes \(p\equiv3\pmod8\) below \(2000\) of good reduction; the valuations of \hyperref[lemma-hlp]{Lemma~9}(c) on a box of
parameters; the reduction types at \(5\) and the degrees of the fifth roots of \(q(d)\) for \(d=1,2,3,5,18,48,68,
193\). Every value quoted in the text from the certificate is asserted, not only printed. No result depends on the
generalized Riemann hypothesis. Apart from the complete enumeration of Cremona's tables in Part~H, which is used only
for \hyperref[thm-25cond]{Theorem~19}, no search is used.

\begin{thebibliography}{LMFDB}

\bibitem[BD]{BD}
N. Bruin and K. Doerksen, \emph{The arithmetic of genus two curves with \((4,4)\)-split Jacobians}, Canad. J.
Math. \textbf{63} (2011): Section~2 (Lemmas~2.4 and~2.6, Proposition~2.8).
\href{https://arxiv.org/abs/0902.3480}{arXiv:0902.3480}.

\bibitem[BK]{BK}
A. Brumer and K. Kramer, \emph{The conductor of an abelian variety}, Compositio Math. \textbf{92} (1994),
227--248.

\bibitem[BL]{BL}
C. Birkenhake and H. Lange, \emph{Complex Abelian Varieties}, 2nd edition, Springer, Berlin, 2004 (Grundlehren der
mathematischen Wissenschaften~302): Sections~5.3 and~12.1 (Corollaries~12.1.2, 12.1.4 and~12.1.5,
Proposition~12.1.3, Lemma~12.1.6) and Corollary~11.8.2.
\href{https://doi.org/10.1007/978-3-662-06307-1}{doi:10.1007/978-3-662-06307-1}.

\bibitem[Cr]{Cr}
J. E. Cremona, \emph{Algorithms for Modular Elliptic Curves}, 2nd edition, Cambridge University Press, 1997; the
tables of elliptic curves over \(\Q\), extended by Cremona to conductor \(500000\), as distributed with Magma.

\bibitem[DS]{DS}
F. Diamond and J. Shurman, \emph{A First Course in Modular Forms}, Graduate Texts in Mathematics~228, Springer, New
York, 2005: Exercise~3.1.1.

\bibitem[Fal]{Fal}
G. Faltings, \emph{Endlichkeitss\"atze f\"ur abelsche Variet\"aten \"uber Zahlk\"orpern}, Invent. Math.
\textbf{73} (1983), 349--366: \S5, Satz~3, Satz~4 and Korollar~1.

\bibitem[HLP]{HLP}
E. W. Howe, F. Lepr\'evost and B. Poonen, \emph{Large torsion subgroups of split Jacobians of curves of genus two or
three}, Forum Math. \textbf{12} (2000), 315--364: Tables~3, 6 and~8.
\href{https://arxiv.org/abs/math/9809210}{arXiv:math/9809210}.

\bibitem[HS]{HS}
M. Hindry and J. H. Silverman, \emph{Diophantine Geometry: An Introduction}, Graduate Texts in Mathematics~201,
Springer, New York, 2000: Theorems~C.1.4 and~C.2.6 and Proposition~C.2.5.

\bibitem[Ka]{Ka}
S. Kamienny, \emph{Torsion points on elliptic curves and \(q\)-coefficients of modular forms}, Invent. Math.
\textbf{109} (1992), 221--229.
\href{https://doi.org/10.1007/BF01232025}{doi:10.1007/BF01232025}.

\bibitem[KM]{KM}
M. A. Kenku and F. Momose, \emph{Torsion points on elliptic curves defined over quadratic fields}, Nagoya Math. J.
\textbf{109} (1988), 125--149: Theorem~(0.1) and the definition of \(Y_1(m,N)\) on p.~125.
\href{https://doi.org/10.1017/S0027763000002816}{doi:10.1017/S0027763000002816}.

\bibitem[Kn]{Kn}
A. W. Knapp, \emph{Elliptic Curves}, Mathematical Notes~40, Princeton University Press, 1992: Theorem~4.2.
\href{https://doi.org/10.1515/9780691186900}{doi:10.1515/9780691186900}.

\bibitem[Ku]{Ku}
D. S. Kubert, \emph{Universal bounds on the torsion of elliptic curves}, Proc. London Math. Soc. (3) \textbf{33}
(1976), 193--237: Table~3.
\href{https://doi.org/10.1112/plms/s3-33.2.193}{doi:10.1112/plms/s3-33.2.193}.

\bibitem[LMFDB]{LMFDB}
The LMFDB Collaboration, \emph{Completeness of elliptic curve data over \(\Q\)},
\href{https://www.lmfdb.org/knowledge/show/rcs.cande.ec.q}{https://www.lmfdb.org/knowledge/\allowbreak show/\allowbreak rcs.cande.ec.q}, 2026.

\bibitem[Ma]{Ma}
B. Mazur, \emph{Modular curves and the Eisenstein ideal}, Publ. Math. IH\'ES \textbf{47} (1977), 33--186:
Theorem~(8). \href{https://www.numdam.org/item/PMIHES_1977__47__33_0/}{Numdam}.

\bibitem[Mo]{Mo}
P. Morton, \emph{Product formulas for the \(5\)-division points on the Tate normal form and the Rogers--Ramanujan
continued fraction}, arXiv:1612.06268. The formulas used here are verified directly by the certificate.

\bibitem[Na]{Na}
F. Najman, \emph{Complete classification of torsion of elliptic curves over quadratic cyclotomic fields}, J. Number
Theory \textbf{130} (2010), 1964--1968: Section~1.
\href{https://doi.org/10.1016/j.jnt.2009.12.008}{doi:10.1016/j.jnt.2009.12.008}.

\bibitem[Ol]{Ol}
L. D. Olson, \emph{Points of finite order on elliptic curves with complex multiplication}, Manuscripta Math.
\textbf{14} (1974), 195--205. \href{https://doi.org/10.1007/BF01171442}{doi:10.1007/BF01171442}.

\bibitem[P256]{P256}
Census proof of 11 November 2026, \emph{Six exclusions at rational torsion order \(256\)},
\href{https://genus-2-torsion.github.io/proofs/qsplit-order-256-obstruction.html}{https://genus-2-torsion.github.io/\allowbreak proofs/\allowbreak qsplit-order-256-obstruction.html}: Section~2.

\bibitem[Se]{Se}
J.-P. Serre, \emph{Propri\'et\'es galoisiennes des points d'ordre fini des courbes elliptiques}, Invent. Math.
\textbf{15} (1972), 259--331: \S1.11 (Propositions~11 and~12 and the Corollaries) and \S1.12 (Proposition~13 and its
Corollary).

\bibitem[Si1]{Si1}
J. H. Silverman, \emph{The Arithmetic of Elliptic Curves}, 2nd edition, Graduate Texts in Mathematics~106,
Springer, New York, 2009: Propositions~IV.3.2, VII.2.1, VII.2.2, VII.3.1, VII.5.1, VII.5.4, VII.5.5 and~VIII.1.5,
Theorems~IV.6.1, IV.6.4, V.2.3.1, V.3.1, VII.6.1 and~VII.7.1, Corollaries~VII.7.2 and~VII.7.3, Remark~VII.1.1 and
Example~V.4.5.

\bibitem[Si2]{Si2}
J. H. Silverman, \emph{Advanced Topics in the Arithmetic of Elliptic Curves}, Graduate Texts in Mathematics~151,
Springer, New York, 1994: IV.9.4 (Tate's algorithm, Step~2), Remark~IV.9.2.2, IV~\S10 and Theorem~IV.10.2,
Theorems~V.3.1 and~V.5.3, Lemma~V.5.2, Corollary~V.5.4 and Proposition~V.6.1.

\bibitem[ST]{ST}
J.-P. Serre and J. Tate, \emph{Good reduction of abelian varieties}, Ann. of Math. (2) \textbf{88} (1968),
492--517: \S1, Theorem~1 and Corollary~2, and \S2. \href{https://doi.org/10.2307/1970722}{doi:10.2307/1970722}.

\end{thebibliography}

\end{document}
